% ============================================================
%  Harald Høffding, "Den nyere Filosofis Historie" (1894–95; 2nd ed. 1903)
%  ENGLISH TRANSLATION — A SELECTION: the Introduction and the Tenth Book
%
%  WHY A SELECTION, AND WHY THESE PARTS.
%  The whole work is ten Books over two volumes, ~2.29 million characters of
%  Danish, and it is already translated: B. E. Meyer, "A History of Modern
%  Philosophy" (Macmillan, 1900), 2 vols. A full retranslation would duplicate
%  Meyer to no purpose. But Meyer has two limits that make a selection worth
%  making:
%
%    1. Meyer's title page says "translated from the German edition" — that is,
%       from Høffding's German of 1895–96, itself made from the FIRST Danish
%       edition. The English therefore stands at two removes from the Danish,
%       and renders an earlier state of the text.
%    2. The edition used here is the SECOND Danish edition (© 1903). Its
%       revisions are concentrated in the notes — some thirty citations dated
%       1900–1903 — but the Introduction is substantively later than Meyer:
%       it refers the reader to "Filosofiske Problemer (Universitetsprogram
%       1902)", a work published two years AFTER Meyer appeared. That
%       cross-reference cannot be in the English.
%
%  The Tenth Book is taken for its subject rather than its date: Robert Mayer
%  and the conservation of energy, then materialism, Lotze, Fechner, von
%  Hartmann, Lange and Dühring. It is the historical backing for the energy
%  argument in Filosofiske Problemer I §4 and in Totalitet som Kategori, both
%  already translated in this edition.
%
%  SOURCE. ~/bibliotek/Høffding, Harald/Den_nyere_filosofis_historie.epub
%  (SAGA / Lindhardt og Ringhof, © 1903, 2020; ISBN 9788726363043), DRM-free.
%  Introduction = OPS/s004-Introduction-01.xhtml (no notes).
%  Tenth Book   = OPS/s015-Chapter-010.xhtml, notes 218–236 in
%  OPS/s017-Notes-01.xhtml, keyed rw-num-note-N.
%
%  PAGE MAP — THERE IS ONE. An earlier survey of this file reported that the
%  endnotes carry no printed-page references. That was wrong: they do, but in
%  the form "218 S. 495." rather than the "(S. 495)" that Filosofiske Problemer
%  uses, and a regex requiring the parentheses found nothing. 231 of the 236
%  notes carry a page, and they map the whole work over pp. 4–565. Volume II
%  restarts its pagination at note 101 (S. 4); the TENTH BOOK runs from about
%  p. 493 to p. 568 of volume II, pinned by notes 218 (p. 495) through 236
%  (p. 565).
%
%  Page markers below are exact where a note pins them and are not interpolated
%  in between; a % [p. N] comment therefore appears at note anchors only.
%  A few of the work's page refs are out of sequence (notes 86, 153) and may be
%  mis-set, as note 43 is in Filosofiske Problemer — neither falls in the Tenth
%  Book.
%
%  DEDICATION. The work as a whole is dedicated "TIL MINDET OM MIN LÆRER OG VEN
%  PROFESSOR HANS BRØCHNER" — to the memory of Hans Brøchner, whose own
%  Problemet om Tro og Viden is transcribed and translated elsewhere in this
%  edition, and whose 1873–74 history Høffding names in the Introduction as
%  the predecessor he means to continue.
% ============================================================
\documentclass[12pt,a4paper]{book}
\usepackage[utf8]{inputenc}
\usepackage[T1]{fontenc}
\usepackage{libertinus}
\usepackage{libertinust1math}
\usepackage{textalpha}    % α, β, γ head the subsections of the Tenth Book
\usepackage[english]{babel}
\usepackage[protrusion=true,expansion=true]{microtype}
\usepackage{setspace}
\usepackage[margin=1.2in]{geometry}
\usepackage{fancyhdr}
\setlength{\headheight}{15pt}
\usepackage{enumitem}
\usepackage{hyperref}

\hypersetup{
  pdftitle={A History of Modern Philosophy: Introduction and Tenth Book},
  pdfauthor={Harald Høffding},
  hidelinks
}

\pagestyle{fancy}
\fancyhf{}
\fancyhead[LE,RO]{\thepage}
\fancyhead[RE]{\textit{A History of Modern Philosophy}}
\fancyhead[LO]{\textit{\leftmark}}

\setstretch{1.3}

\title{\textbf{A History of Modern Philosophy}\\[6pt]
  \large Introduction and Tenth Book}
\author{Harald Høffding\\[6pt]
  \normalsize translated from \textit{Den nyere Filosofis Historie},
  2nd edition (1903)}
\date{1894--1895 / 1903}

\newcommand{\hhchapter}[1]{%
  \chapter*{#1}%
  \addcontentsline{toc}{chapter}{#1}%
  \markboth{#1}{#1}%
}

\begin{document}
\maketitle
\thispagestyle{empty}
\newpage

{\small\noindent\textit{Translator's note.} This volume contains two parts of
Høffding's two-volume history: the Introduction, in which he sets out his
conception of philosophy and of what a history of philosophy is for, and the
Tenth Book, on the philosophy of the second half of the nineteenth century.
The remainder of the work is available in English in B.~E. Meyer's translation
of 1900, made from the German edition. The Introduction translated here is
later than Meyer's text and refers to a work of 1902; the Tenth Book is given
because its subject --- Robert Mayer, the conservation of energy, and the
idealist and critical reactions to materialism --- is the background to
arguments translated elsewhere in this edition.\par}

\newpage
\tableofcontents
\newpage

% ============================================================
%  INDLEDNING  (OPS/s004-Introduction-01.xhtml; 11 paragraphs, no notes)
% ============================================================
\hhchapter{Introduction}
\label{ch:indledning}

It will be natural that, when a new presentation of an important period in the
history of philosophy appears, the question should be asked what conception of
philosophy the author takes as his ground, and what significance and value he
ascribes to the history of philosophy. My answer is that the studies whose
provisional conclusion is marked by the publication of this book had precisely
the aim of orienting me with respect to what philosophy really is. As we come to
know a person through his biography, so we must be able to come to know a
science through its history. And it will be all the more natural to try this
road, since experience constantly shows that in the domain of philosophy
conflicting ways of seeing assert themselves, so that there is no presentation
of philosophy to which one could point as giving an exhaustive concept of it. It
will then be natural here --- as in the domain of religion --- to apply the
comparative method. The history of philosophy treats the attempts which
individual thinkers have made to discuss the ultimate problems of knowledge and
of life. The most important result will be to get clear about which problems
philosophy occupies itself with, how those problems are posed at different
times, and by what the different ways of posing them, and of attempting to solve
them, are determined in detail. If such a result can be approximately reached,
the study of the history of philosophy will be of no slight importance for
continued philosophizing.

For my own part I have been strengthened, by the study of the history of modern
philosophy, in the assumption I had arrived at by another road: that
philosophical inquiry turns upon four principal problems. These four problems
are now to be briefly characterized. A more searching discussion and comparison
of them I have given in the treatise \textit{Filosofiske Problemer} (University
Programme 1902).

\medskip

\noindent\textbf{1. The problem of knowledge (the logical problem).} However
different the several sciences may be in respect of subject and of method, they
all nevertheless work with human thinking. Every time they form a concept, utter
a judgment or draw an inference, they presuppose the general forms and
principles of thinking. Here, then, there shows itself the possibility of a
peculiar discipline, one which investigates the forms in which thinking moves
and the principles it must by its nature take as its ground, whatever the
objects it may concern. This discipline, \emph{formal logic}, treats however
only one part of the problem of knowledge. Those forms and principles do not
lead beyond thinking itself, but only make it possible for thinking to be in
agreement with itself, to be consistent. Every time they are applied to given
phenomena, which our thinking has not itself constructed but must take as they
come, the question arises with what right that application takes place ---
with what right we assume that not only our thinking, but existence itself,
which expresses itself in the given phenomena, is consistent, agrees with
itself. Herewith shows itself the possibility of a discipline which investigates
the conditions of a knowledge of existence, and the limits of such knowledge.
This discipline is \emph{epistemology}.

\noindent\textbf{2. The problem of existence (the cosmological problem)} arises
when it is asked how the existence of which we are members must be supposed to
be constituted, when we draw the ultimate consequences of all that we know, or
can conjecture with the help of the most probable hypotheses. We call this
problem the cosmological (from \textit{kosmos}, world, and \textit{logos},
doctrine), because it leads to a discussion of the possibilities that present
themselves to thought when it seeks to work up what is given in experience into
a general world-view, or when by speculation's bold road it attempts to
construct one. The different philosophical systems are attempts in this
direction --- attempts whose value depends on how significant and how
comprehensive are the experiences they take as their ground, and how much
consistency and power of combination is displayed in their constructions.

\noindent\textbf{3. The problem of valuation (the ethical-religious problem)}
arises through this, that we do not merely stand to existence in an intuiting
and understanding relation, but also have our feeling set in motion by it, so
that we utter judgments which ascribe value to it or deny it value. Of
particular importance are the judgments we utter upon human actions, our own and
others'. All such appraisal rests --- like all knowledge and understanding ---
upon certain presuppositions, which it is a matter of exhibiting and
determining. This task belongs under the ethical problem. When the valuation
concerns not merely human actions and institutions but existence, life as a
whole, the religious problem arises, which leads especially to a discussion of
the relation between the ethical ideals and actual existence, so that it becomes
a combination of the cosmological and the ethical problem.

\noindent\textbf{4. The problem of consciousness (the psychological problem).}
For all three of the problems indicated it shows itself that their treatment
presupposes an empirical acquaintance with human conscious life. Psychology
presents the \emph{actual} development of human knowledge, which must be known
before one can discuss the conditions of knowledge's \emph{validity}. And since
the relation between the spiritual and the material is a principal point within
the problem of existence, psychology is presupposed in cosmology as well. As
regards the problem of valuation, finally, psychology discusses partly the
constitution of the very feelings that lead us to make valuations, partly the
possibilities which actual conscious life presents for a further development in
the direction valuation demands. In consequence of this close relation to the
philosophical problems, psychology is itself to be reckoned a part of
philosophy; and one will therefore also, conversely, quite naturally come in
psychology to touch upon philosophical problems. And even if one were to declare
all three of the preceding problems insoluble, or arisen through
misunderstanding, the psychological problem --- the question of the nature and
the laws of conscious life --- would still remain as philosophy's last
stronghold.

\medskip

--- If it is asked what factors will, in the nature of the case, come to
influence the treatment and the answering of these problems, the first to be
named must be the personality of the one who philosophizes. The problems named
have this in common, that they lie at the boundary of our knowledge, where the
exact methods fail; it cannot then be otherwise than that the inquirer's
personality must come to determine his train of thought, without his needing to
be conscious of it. The ``personal equation'' must be of greater importance in
philosophy than in other scientific domains. The historical, comparative method
therefore becomes of particular importance here; by its help the personal
element can most easily be exhibited. That personal element is not always
something to be pushed aside; its presence is often a condition of a problem's
being able to come to be at all. There are thoughts which can come to be only
upon a determinate psychological soil. --- It will, in the second place, depend
upon which observations are taken as the ground. Here the development of natural
science acquires especially great importance for modern philosophy. As will
appear, the decisive problems of modern philosophy are determined by the very
fact that modern natural science came into being. But to this there come as well
--- especially as regards the problem of valuation --- historical conditions,
intellectual movements in other domains. --- Finally, the posing of problems and
their solution are determined by the consistency with which the starting points
once seized are held fast and carried through.

Both the characterization and the criticism of every philosophical attempt will
always lead back to these three points. Among them I lay most weight, in the
presentation that follows, upon the first two. An inconsistency in a great
thinker is often precisely a natural consequence of \emph{several} trains of
thought having opened themselves to his genius, without his having been able to
pursue them far enough to discover their mutual conflict. And it may be of the
greatest importance that these different trains of thought be brought out. Best
of all, of course, is when profundity and consistency are united.

An attempt at posing or at solving a problem may have a twofold interest. In the
first place it may be regarded as a symptom, as the historical expression of
intellectual currents. From this side the history of philosophy is a part of
general cultural history. In the second place it may be examined with respect to
the real, definitive clarity that has been reached by it. This consideration is
more specifically philosophical and technical. The two sides from which the
philosophical phenomena can be regarded will of course, in the case of different
phenomena, stand in very different relations to one another. Sometimes the
cultural-historical interest may even stand in inverse relation to the purely
philosophical.

If I am briefly to indicate wherein my presentation of the history of modern
philosophy differs from its predecessor in our literature --- Brøchner's work of
1873--74 --- I would name first the particular weight I lay upon the personal
factor, upon the relation to empirical science, and upon the cultural-historical
significance of the philosophical phenomena. Next I would point to my endeavour
to give the posing of problems precedence over their solution. Solutions may
die, and the problems live on; otherwise philosophy would not have had so long a
life as it has had. Finally it goes without saying that a renewed study of the
sources, as well as the use of the abundant literature of recent years, has set
many phenomena in a new light. But in spite of all this it has been a dear
thought to me, in this work, to be able to continue Brøchner's labour in the
field of the history of philosophy; and I can only hope that my book may not be
unworthy to be consecrated to the memory of that noble inquirer, whose
admiration for the heroes of thought was at once so deep and so understanding.

\begin{center}---\end{center}

% ============================================================
%  TIENDE BOG  (OPS/s015-Chapter-010.xhtml; 153,793 chars; notes 218--236)
%
%  Section map, by character offset within the Book. Offsets come from the
%  cumulative paragraph lengths of the epub file and are the batching key,
%  since there is no pagination to use:
%
%     0        TIENDE BOG — opening
%     3,465    1. Robert Mayer og Principet om Energiens Bestaaen
%    15,203    2. Materialisme
%    33,947    3. Idealistiske Konstruktioner paa realistisk Basis
%    33,947       a. Rudolph Hermann Lotze
%    42,449          α. Den mekaniske Naturopfattelse
%    51,346          β. Metafysisk Idealisme
%    61,061          γ. Spiritualistisk Psykologi
%    70,747       [b.] Gustav Theodor Fechner        <- printed "2."; see below
%    74,762          α. Poetisk spekulativ Verdensanskuelse
%    80,015          β. Psykofysik
%    86,374          γ. Naturfilosofi
%    89,058       c. Eduard von Hartmann
%    91,570          α. Naturfilosofi og Psykologi
%    99,362          β. Pessimisme og Etik
%   106,471    4. Kriticisme og Positivisme
%   106,471       [a.] Friedrich Albert Lange        <- printed "α."; see below
%   127,375       b. Eugen Dühring
%   135,522          α. Erkendelsesteori
%   143,694          β. Verdensopfattelse
%   150,192          γ. Etik
%
%  TWO NUMBERING FAULTS IN THE SOURCE, to be reproduced as printed and logged,
%  not silently mended:
%    · under §3 the sub-heads run a. / 2. / c. — Fechner is set "2." where the
%      series requires "b."
%    · under §4 the sub-heads run α. / b. — Lange is set with the Greek letter
%      belonging to the level below, where the series requires "a."
%  The epub also alternates β and ß in the Greek sub-heads; ß is a mis-set
%  eszett for beta and is normalised to β here.
% ============================================================
\hhchapter{Tenth Book}
\label{ch:tiende-bog}

\begin{center}\textsc{Philosophy in Germany 1850--1880.}\end{center}

\bigskip

\noindent Of the two great tendencies in nineteenth-century philosophy it was
positivism that had best maintained its connection both with the thought of the
eighteenth century and with empirical science. The philosophy of Romanticism, by
contrast, was a definite and conscious attempt at reaction against both these
tendencies; it wanted even to do over again what had been founded with the rise
of natural science in the seventeenth century. In Germany, the hearth of
Romanticism and of its philosophy, that tendency was the dominant one at the
middle of the century. The adherents of the critical undercurrent were the only
ones who held fast to continuity with the other sciences --- leaving aside
Schopenhauer and Feuerbach, who at that time still stood as solitary thinkers
for whom the hour of recognition had not come.

As little as positivism was to be explained as a reaction against Romanticism,
so little is the philosophical movement which takes place in Germany after the
middle of the century to be explained as a continuation of Franco-English
positivism. It has its presuppositions in Germany itself. It receives its
problems above all from the renewed rise of natural science towards the middle
of the century. Not only do natural-scientific studies and results now become
the object of more general interest than in the earlier part of the century,
when poetry, religion and speculative philosophy had laid claim to men's minds;
natural science itself now returned, with clearer consciousness, to the great
principles its founders had in their day set up. The demand for an
understanding of outer nature as a series of demonstrable causes and effects ---
in other words, for a mechanical explanation of nature --- was now asserted
again, as in the seventeenth century. And at the same time a new and
comprehensive law of nature was discovered and demonstrated: the law that in
physical nature no energy arises or perishes, but that what happens in every
apparent arising or perishing is a conversion of energy into another form, and
in such a way that the different forms of energy stand in a determinate
quantitative relation to one another. This law --- which, together with Darwin's
doctrine of the origin of species through the struggle for existence, is the
most significant result of natural-scientific investigation in the nineteenth
century --- could not but set philosophical thinking in motion, just as in their
day the Copernican theory and Galilei's founding of mechanics had led philosophy
into new paths. For German philosophy in particular it had to become a principal
question: how far can the thoughts developed by the philosophy of Romanticism be
held fast, if one accepts the new natural-scientific way of regarding things?
This question was answered in three different ways. Modern materialism rejected
those thoughts as pure illusions and proclaimed the doctrine that only the
material exists, as a necessary consequence of natural science. Lotze and
Fechner, by contrast, attempted to exhibit the fundamental thought of speculative
philosophy of religion as the last and concluding presupposition of the picture
of existence that can be built up by the help of the natural-scientific method.
Finally Albert Lange and Eugen Dühring approached critical philosophy and
positivism, by bringing out the significance of the problem of knowledge and by
maintaining practical idealism's independence of empirical science, while at the
same time insisting upon experience's right to determine the actual content of
the world-picture.

\section*{1.\quad Robert Mayer and the Principle of the Conservation of Energy}
\label{sec:mayer}

Not only in philosophy, but in the other sciences too, knowledge moves forward
in rhythmic fashion, through actions and reactions. At the close of the
eighteenth century natural science had a great and fruitful period, in which
some of the most important truths in the domains of chemistry and physiology
were discovered. \textsc{Lavoisier} introduced the quantitative method into
chemistry, and by its help the truth of the old thought was established that no
matter arises or perishes, but that the same mass of matter subsists through all
changes, under different forms. Chemistry was founded as an exact science.
\textsc{Priestley}, \textsc{Ingenhouss}, \textsc{Senebier} and \textsc{Saussure}
found the principal laws of the metabolism of plants and animals, and thereby
founded the great doctrine of the circulation of matter in nature, whereby the
organic and the inorganic world appear in their close mutual connection. It was
discovered that it is under the influence of sunlight that organic matter is
formed in the green plant cells by the taking up and dissolution of the carbonic
acid of the air. The carbon stored in the plant cells then serves animals for
nutriment; by the animal functions a combustion takes place, and the carbonic
acid thereby formed passes out into the air in respiration, whence the same
circulation can begin anew. A great cosmic connection had here been
% [p. 495]
exhibited.\footnote{On ``the period of the great discoveries'' at the close of
the eighteenth century see Rasmus Pedersen: \textit{Planternes Næringsstoffer.}
Historisk Indledning. København 1883, p.~44\,ff.} The picture of the world
sketched by Copernicus, Bruno, Kepler and Galilei, which Newton had enriched by
exhibiting universal attraction as a binding and ordering tie, hereby received a
new and significant enrichment.

But the age was too much taken up with revolution and world war, with
Romanticism and speculation, and soon with orthodoxy and mysticism as well, for
these great ideas to gain the influence upon general opinion that they deserved.
Even within natural science their victory was hindered by other tendencies. It
was within the science of organic life that resistance was chiefly offered to
the new points of view, because men were still governed by the reaction against
the Cartesian doctrine of the organism as a machine. The expressions peculiar to
the organism were derived from a peculiar vital force, supposed to be wholly
different from the other forces of nature. This mode of explanation, so-called
vitalism, could not but lead to a rejection of the attempts to draw the organic
phenomena into the universal circulation of
% [p. 496]
matter.\footnote{On vitalism in its various forms see Louis Peisse: \textit{La
Médicine et les Médicins.} Paris 1857. I, p.~226--297. --- It must be firmly
held that by vitalism is here meant a tendency which finds the
\emph{explanation} of vital phenomena in the assumption of a ``vital force'' ---
not a tendency which maintains that there is much in the vital phenomena that
has hitherto not been explained by physical and chemical causes. See my treatise
\textit{Om Vitalisme} (in ``Mindre Arbejder'').} Besides, the prevailing
interest among botanists and zoologists was to describe the forms and determine
their systematic place, not to find the developmental process through which the
forms come to be, and the causes that determine that process. Romantic
philosophy of nature favoured this aesthetic and formal conception --- and had
itself proceeded from it. When ``philosophy of nature'' is often given the whole
blame for the fact that the points of view of the close of the eighteenth
century did not sooner prevail in botany and zoology, it must not be forgotten
that it was quite as much an effect of the condition of natural science as a
cause of it.

From about 1840 a change came about in this respect. Life began to be conceived
as something other than a play of forms or a revelation of ideas.
\textsc{Dumas} and \textsc{Liebig} insisted upon the significance of the
chemical processes for plant and animal life. And from the medical side too a
strict mechanical explanation began to be demanded of the processes that go on
in the organism. In some of his earliest writings \textsc{Hermann Lotze} worked
in this direction, especially in \textit{Allgemeine Pathologie und Therapie als
mechanische Naturwissenschaften} (1842) and in the treatise \textit{Leben,
Lebenskraft} (1843, in R. Wagner's Handwörterbuch der Physiologie). The task was
set of making physiology a purely mechanical science. Most epoch-making,
however, was the fundamental proposition set up by the physician and physicist
\textsc{Robert Mayer} (b. 1814, d. 1878) in his work \textit{Die organische
Bewegung in ihrem Zusammenhange mit dem Stoffwechsel} (1845): that in the
process of life there takes place only a transformation of force as well as of
matter, never a creation of them. Mayer here points back to the great law he had
discovered a few years earlier, when he set the law of the conservation of force
beside the law of the conservation of matter deriving from Lavoisier.

Robert Mayer's life as an investigator is borne by a single thought, which he
came upon early; but it was the tragedy of his life that he was unable to carry
it through in detail by experimental means. As a child he had tried to construct
a \textit{perpetuum mobile}, and it made a strong impression upon him that he
did not succeed. From that time he brooded constantly upon the relation between
cause and effect --- first and foremost within physiology, which as a medical
man lay nearest to him, but thereafter within chemistry and physics as well. On
a voyage to the East Indies as ship's doctor he arrived at the new idea of the
imperishability of force in nature, under the influence partly of investigations
into the origin of organic heat, partly of the fact that heat arises in the
wave-motion of the sea. His point of departure is thus definitely observation.
But that observation led him to ask whether a law similar to the one in
chemistry might not hold in physics too, so that not only matter but force also
was imperishable, and every change a mere conversion. And he then came back to
his earlier broodings upon cause and effect. It was for him a self-evident
proposition that there could not be more in the effect than in the cause:
\textit{causa æqvat effectum}! It follows from ``the law of the logical
ground''! Mayer thus does not distinguish, any more than the dogmatic
philosophers do, between ground and cause; thereby he comes, like William
Hamilton and Herbert Spencer, to regard it as a matter of course that the cause
cannot become nothing when the effect arises, but that the effect must furnish
an equivalent for what seemed to vanish with the cause. Hitherto men had rested
content with the fact that a motion ceased when it met a sufficiently strong
resistance, and that heat arose from frictional resistance. But --- Mayer asked
--- does the motion then become nothing, and does heat then arise out of
nothing? If that is so, then ``the red thread of science'' is cut through, just
as it would be if in chemistry we were to assume that oxygen and hydrogen became
nothing when they combined, and that water then came to be out of nothing. Just
as we assume that oxygen and hydrogen are transformed into water, so we must
also assume that the motion does not become nothing, but is transformed into
% [p. 497]
heat.\footnote{Cf.\ Mayer's letter to Griesinger of 20 July 1844:
``\textit{Wenn ich sage, Wärme lässt sich in Bewegung verwandeln, und umgekehrt,
so will dies nichts meinen, als dass zwischen Wärme und Bewegung finden hin und
her dieselben qvantitativen Beziehungen statt.}'' (\textit{Kleine Schriften und
Briefe} von Robert Mayer. Herausg.\ von Weyrauch. Stuttgart 1893, p.~225.) ---
In his paper \textit{Robert Mayers Entdeckung und Beweis des Energieprincipes}
(in the Festschrift for Sigwart. Tübingen 1900) A. Riehl illuminates Mayer's
doctrine from an epistemological point of view.} And by the expression
\emph{transformed} Mayer meant to designate only that a constant quantitative
relation subsists between the cause, which disappears, and the effect, which
comes in its place. If the train of thought put forward is correct, experience
must establish such a relation of equivalence between the different forms of
force that succeed one another. On the basis of experiments already available,
Mayer sought already in his first paper (\textit{Bemerkungen über die Kräfte der
unbelebten Natur}, printed in Liebig's ``Annalen der Chemie und Pharmacie'',
1842) to determine the quantitative relation between heat and motion. It was
this idea of an invariable relation of magnitude that was the main thing for
Mayer, and that in the following period proved so fruitful for natural science.
It was a significant extension of the abstract law of the conservation of force,
already maintained by Huyghens and Leibniz, so as to hold for the relation of
the different forces of nature to one another. Mayer derived from his train of
thought that there really exists only a single force, which appears under
different forms, and these stand to one another in a determinate quantitative
relation.

Although Mayer's attempt to give a philosophical derivation of the law of the
conservation of energy does not hold, his train of thought has nevertheless an
epistemological interest. For it was in itself fully legitimate to raise the
question whether between the event we call the cause and the event we call the
effect there might not obtain a relation similar to that between logical ground
and logical consequent. And Mayer's train of thought leads naturally to raising
that question. Experience must then decide what the answer is to be. --- It is
interesting that the Danish physicist \textsc{Colding}, who in the year after
Mayer's first paper was printed set up the proposition of the conservation of
energy by a road of his own and confirmed it by experiment, proceeded from the
view that it was really a proposition of reason; and likewise that
\textsc{Helmholtz}, in his paper \textit{Ueber die Erhaltung der Kraft} (1847),
proceeded from an epistemological postulate. The Englishman \textsc{Joule}, who
like Colding and Helmholtz had come by a road of his own to the same result as
Mayer --- a remarkable example of how several may simultaneously be on the track
of one and the same discovery --- proceeds predominantly experimentally, but
nevertheless indicates the a priori improbability that force should be
% [p. 499]
annihilated without a corresponding effect.\footnote{``If the resistance to
electrolysis which is over and above that due to chemical change were not
accounted for elsewhere, \textit{it would prove the annihilation of a part of the
power of the circuit, without any corresponding effect.} We shall see that this
is not the case, but that in the evolution of heat, where the excess of
resistance takes place, an exact equivalent is restored'' (Joule: \textit{On the
Heat evolved during the Electrolysis of Water.} 1843. Scientific Papers. I.
London 1884, p.~115). ``\textit{We might reason, a priori, that absolute
destruction of living force cannot possibly take place,} because it is manifestly
absurd to suppose that the powers with which God has endowed matter can be
destroyed~... but we are not left with this argument alone, decisive as it must
be to every unprejudiced mind'' (Joule: \textit{Matter, Living Force and Heat.}
1847. Scientif.\ Pap.\ I, p.~269). --- Joule's standpoint is therefore not in
principle different from Mayer's and Colding's. For the latter compare an
utterance in his \textit{Undersøgelse om de almindelige Naturkræfter og deres
gensidige Afhængighed.} (Videnskabernes Selskabs Skrifter. Femte Række. Naturv.\
mathem.\ Afd.\ II, p.~129): ``The thought that an activity could disappear in
the corporeal without appearing again as operative cause seems to me contrary to
reason.'' --- It is of interest in this connection to recall that Lavoisier too
proceeded in his day from a fundamental proposition he regarded as a matter of
course, namely that the weight of a compound body is equal to the weight of the
substances of which it consists. Self-evident as that presupposition was, it
nevertheless had to be proved to those who took heat for a substance; for if
that assumption were correct, the release of heat which accompanies a chemical
combination would have to entail a diminution of weight. Compare Ernst von
Meyer: \textit{Geschichte der Chemie.} Leipzig 1889, p.~137.}

The new proposition had difficulty in winning recognition, but gradually made
its way in the course of eighteen years (1842--1860), especially as it proved
fruitful in leading to new investigations and discoveries. Its great importance
consists precisely in the determinate questions it leads the investigator of
nature to ask, every time a peculiar manifestation of force appears or ceases.
Just as it is the importance of the general causal proposition that at every
change which occurs it leads us to look about for a preceding change of which
the new one might be the effect, so it is the importance of the proposition of
the conservation of force --- or, as one now prefers to say, following the
proposal of English investigators, of energy --- that one is at once brought to
set oneself the task of finding out in what relation the manifestations of force
that succeed one another stand to one another.

Philosophically regarded, the question had above all to become how the spiritual
phenomena stood to the new law. It is characteristic in this respect that almost
all its discoverers proceeded from decidedly spiritualistic and teleological
views. Mayer repeatedly declared himself a definite opponent of materialism, and
expressed as his conviction that the scientific truths stand to the Christian
religion as brooks and rivers to the ocean. At the congress of natural
scientists in Innsbruck in 1869 he made a declaration to this effect, to the
offence of Carl Vogt and his friends. Colding's standpoint is seen from the
following utterance: ``My first thought that the forces of nature must be
imperishable I drew from the view that the forces in nature must be akin to the
spiritual in nature, to the eternal reason as well as to the human mind. It was
thus the religious conception of life that led me to the thought of the
imperishability of the forces of nature.'' Colding and Joule conceive that God
at the creation laid a certain sum of force into nature, and that this sum can
neither be increased nor diminished as a whole, but only distributed in
different ways --- as Descartes in his day assumed of motion. And both assumed
that with the conservation of energy the conservation of what is valuable in the
world was secured, since they no more than Leibniz distinguished between energy
% [p. 500]
itself and its employment in favour of life and development.\footnote{Mayer:
\textit{Kleinere Schriften und Briefe}, p.~339\,f., 460. \textit{Die Mechanik
der Wärme.} 3.~Ausg., p.~356\,f. --- Colding: \textit{Naturvidenskabelige
Betragtninger over Slægtskabet mellem det aandelige Livs Virksomheder og de
almindelige Naturkræfter.} (Oversigt over det kgl.\ danske Videnskabernes
Selskabs Forhandlinger. 1856), p.~155, 166. --- Joule: \textit{Scientific
Papers.} I, p.~269, 273. --- As regards Mayer, it appears from his address at
Innsbruck (\textit{Die Mech.\ d.\ Wärme}, p.~357) that he did not assume any
break in the series of material phenomena, but assumed a parallelism between the
spiritual processes and the brain processes.}

But a closer investigation and grounding of these presuppositions the natural
scientists named do not undertake; they proceeded, like the philosophers of the
seventeenth century, from given religious representations. A closer
investigation of the relation of the law exhibited --- and of the knowledge of
nature so far won generally --- to the phenomena of spiritual life had therefore
naturally to be called forth. And it is then no wonder that the ways parted, and
that different tendencies stepped forth against one another. All that was given
was this: that when one kind of physical energy ceases, another kind of physical
energy enters, in a certain determinate quantity, in its place. The question was
what significance this knowledge must consistently come to have for the
conception of the world.

\section*{2.\quad Materialism}
\label{sec:materialisme}

The strongest reaction against the philosophy of Romanticism --- apart from
crass orthodoxy --- is marked by the materialist literature which flourished in
Germany at the middle of the century. If there the Idea had been everything, now
matter became everything. Enthusiasm for natural science, and for the great
points of view it had succeeded in carrying through in the doctrine of the
conservation of matter and of force, could naturally lead to regarding that
doctrine itself as a sufficient philosophy, and to explaining every side of
existence by its help. The simplest way of solving the question what
significance the law of the conservation of matter and force has for our
world-view would undeniably be to say that this law contains our whole
world-view. The philosophy which modern materialism sets up will be nothing
other than the simple consequence of natural science. Thereby a firm basis
seemed at the same time to have been won for all our representations and for our
conduct of life. In place of a mystical or spiritualistic foundation we now get
a wholly tangible foundation to build upon in theory and in practice. And it is
a doctrine easy to popularize. It is childlike, intuitable, easily accessible,
and its presentation gives good occasion to bring forward a mass of
natural-scientific facts of general interest. Several of the spokesmen of German
materialism were able and independent investigators of nature, such as
\textsc{Carl Vogt} and \textsc{Jacob Moleschott}. Others, like \textsc{Louis
Büchner}, worked especially by clear, tasteful and enthusiastic presentation. It
is in its popularizing tendency that this literature has its greatest merit. A
wealth of knowledge has been spread abroad by it in wide circles. Büchner's
\textit{Kraft und Stoff} is one of the most read popular-scientific books of our
century; from 1855 to 1899 nineteen German editions appeared, and it has besides
been translated into several other languages. Even if there is something
dogmatic about materialism, it has nevertheless done good service by setting its
dogmatism against ecclesiastical dogmatism, and thereby awakening reflection
upon problems which had been pushed aside at the dissolution of Romantic
philosophy. The whole materialist movement in Germany was moreover borne by an
idealistic interest in humanity and progress, and Büchner was fully justified in
protesting against the confusion of materialism as method and theory with
materialism in the sense of a practical direction of life. Materialism can very
well acknowledge the value of the highest and noblest ideas and feelings ---
even if it holds that they, like all spiritual phenomena, are only products of,
or forms of, material processes.

When materialism claims to be the simple consequence of the results of natural
science, it is of interest to observe --- not only that the discoverers of the
law of the conservation of energy proceeded throughout from spiritualistic
presuppositions, but that the most important collisions in the materialist
controversy took place between parties who all stood upon the ground of natural
science. This does not necessarily show that materialism is wrong; but it shows
how difficult it is to draw the true consequences, and how many different
motives come to influence the world-view of individuals. Moleschott's famous
book \textit{Der Kreislauf des Lebens} (1852) is directed against
\textsc{Liebig}'s theologizing utterances in his ``Chemische Briefe'', in which
he had particularly attacked an earlier proposition of Moleschott's: ``Ohne
Phosphor kein Gedanke!'' At the same time a controversy broke out between the
physiologist \textsc{Rudolph Wagner} in Göttingen and the zoologist \textsc{Carl
Vogt} in Geneva. This controversy reached its height when Wagner, at a congress
of natural scientists in Göttingen in 1854, declared himself in favour of the
assumption of an ethereal soul-substance which moves the brain's fibres as the
musician moves the pianoforte, and which is propagated from parents to children
by division [!] --- an assumption he supported upon the Bible's teaching, but
which he admitted he could maintain only by a sharp distinction between faith
and knowledge. For when one goes beyond the domain of exact knowledge one has,
according to Wagner, only an unprovable faith to hold to --- and he adds: ``In
Sachen des Glaubens liebe ich den schlichten, einfachen Köhlerglauben am
meisten.'' Against this Vogt came forward with a robust polemic:
\textit{Köhlerglaube und Wissenschaft} (1855). In a later work,
\textit{Vorlesungen über den Menschen} (1863), Vogt gives a more scientific
presentation of his standpoint. He lays weight above all upon the brain's being
the organ of consciousness, and upon consciousness standing to the brain as
every function stands to the corresponding organ. Although he admits that it
cannot be explained \emph{how} consciousness arises in the brain's cells, he
holds fast that it is indissolubly bound to them; and he does not give up the
teasing proposition he had already set up earlier, that thought stands in the
same relation to the brain as bile to the liver or urine to the kidneys ---
notwithstanding that a function stands in a different relation to its organ than
a product does to the place where it is produced.

Moleschott relies, as the title of his work shows, especially upon the doctrine
of the conservation of matter. It is for him the great thought which the
Encyclopaedists of the eighteenth century had already brought forward, which
subsequent natural science has confirmed, and upon which the faith of the future
must be built. He feels reverence for the great circulation in nature. The miner
fetches phosphate of lime up out of the earth in the sweat of his brow --- and
perhaps thereby the matter destined for the best head and the highest thoughts
passes through his hands: the farmer manures his soil with the phosphate of
lime; this becomes a constituent of the wheat that nourishes men's bodies and
brains! With matter, life circulates through all parts of the world; with life,
thought --- and from thought there springs in turn the will to make life better
and happier. If we can but supply the organism and the brain with the best
substances, thought and will too will attain their highest development. The
investigator of nature is our age's Prometheus, and chemistry is the highest
science. The social question finds its solution if one can but [sic] find the
right distribution of the substances to which the life of thought and of will is
bound. --- Only the first edition of Moleschott's book has been taken into
account here; in later editions it was altered and expanded. Moleschott, who was
born in Holland in 1822, was working as a lecturer in Heidelberg when the book
appeared; but when his freedom of teaching was interfered with there, he moved
to Zürich; later he worked as professor of physiology in Turin and in Rome,
where he died in 1893. After his death his autobiography appeared
(\textit{Für meine Freunde.} Lebenserinnerungen von Jacob Moleschott. Giessen
1895). It gives a fine picture of the idealistically minded investigator of
nature, striving after an all-round culture. For a right understanding of his
standpoint the following utterance is of service, which comes after the train of
thought of ``Kreislauf des Lebens'' has been developed: ``One-sidedly
`materialistic' this is only for those who can imagine a matter without force or
a force without material bearer. I was clearly conscious that one could turn the
whole way of representing it round: since all matter is a bearer of force, is
endowed with force or penetrated by spirit (\textit{geistesdurchdrungen}), one
may just as well call it an idealistic view'' (p.~221). Moleschott's standpoint
is rather to be called monism than materialism: ``It is a question of a true,
indivisible two-in-oneness, and the dispute is not between materialism and
spiritualism, but between two-in-oneness (\textit{Zweieinigkeit}) and divided
duality (\textit{zwiespältige Anschauung})'' (p.~222).

\textsc{Louis Büchner} (1824--1899) does indeed declare that he will not explain
what the relation is between spirit and matter, force and substance; he will
only maintain that there is a necessary and indissoluble relation between them.
But even this is more than can be scientifically proved; and Büchner moreover
harbours no doubt that spirit is only a property of matter, force only a
property of substance. The impulse to his famous book ``Kraft und Stoff'' he had
received from reading Moleschott's ``Kreislauf des Lebens'', and it is the
eternity of matter that is the ultimate foundation in him too. He holds that the
conservation of force is really a self-evident consequence of the conservation
of matter, and could therefore already have been drawn from Lavoisier's
chemistry. It was only in the fifth edition of his book, however, that he
inserted a special chapter on ``the immortality of force'', in which the
circulation of force is set as a necessary correlate beside the circulation of
matter, so that ``the two together form, from eternity and to eternity, the sum
of phenomena which we call the world''. Yet he constantly holds fast to the
assumption that ``in matter dwell all the forces of nature and of spirit; matter
is the primal ground of all being''. And as regards the spiritual, he regards it
--- notwithstanding that he has disclaimed all insight into its relation to
matter --- as a mere product. ``As the steam engine produces motion, so the
complicated organic composition of force-endowed substances in the animal body
produces a total sum of certain effects which, when they are bound together into
a unity, are called by us spirit, soul, thought'' (Kraft und Stoff. 7.~Aufl.,
p.~130). In a later utterance Büchner goes --- without appearing to notice it
--- a step further. Thus it runs, in an utterance directed polemically against
the author of the present work: ``The opposition between physical and psychical
energy is tenable only so long as one conceives the concepts of body and spirit,
or more generally of matter and force, in a dualistic way, whereas for the
materialistic-monistic conception they coincide into one and perhaps present
only \emph{different sides or forms of expression of one and the same primal
ground of all things}. It is precisely the law of the conservation of force that
in the problem of the soul leads necessarily to materialistic consequences, as
I, supported by physiological-psychological facts and considerations, believe I
have established to evidence in several of my writings, and recently again in
the fifth of my letters on \textit{Das künftige Leben und die moderne
Wissenschaft} (Leipzig 1889). I believe I have there shown that psychical
activity \emph{is not, and cannot be, anything other than the radiation, taking
place between the cells of the grey cerebral cortex, of a movement initiated by
outer impressions}'' (from an article in the weekly \textit{Menschthum.} Gotha
1889, No.~46).

It is clear that Büchner in this utterance confuses two different theories: the
one, which conceives spirit and matter as forms of expression of one and the
same substance, and the other, for which matter is ``the primal ground'' and
psychical activities are therefore consistently \emph{nothing other than} motion
(or ``radiation of motion''). In one and the same breath he professes both. The
tireless champion of the influence of natural-scientific results upon our
world-view has not seen this inconsistency, which perhaps in his eyes did not
much matter, since his real aim was not to inculcate any one determinate
philosophical theory. In his great zeal he has also believed that he had
``established to evidence'' more than can be established.

With sharper criticism \textsc{Heinrich Czolbe} (1819--1873), a physician like
Büchner, sought to determine the consequences of the teaching of natural
science. In his first works he is a consistent materialist. In \textit{Neue
Darstellung des Sensualismus} (1855) and \textit{Die Entstehung des
Selbstbewusstseins} (1856), as well as in the paper \textit{Die Elemente der
Psychologie vom Standpunkte des Materialismus} (Zeitschrift für Philosophie.
XXVI), he develops the view that materialism conflicts with natural science so
long as the latter assumes the doctrine of the specific energy of the sensory
nerves, and thus assumes a difference between the sensations and the outer
processes corresponding to them. In opposition to that doctrine Czolbe sets up
the assumption that it is essentially one and the same motion which propagates
itself from the surrounding world, through sense-organs and nerves, to the
brain. At no point does any qualitative change take place; the difference
between the different senses is to be explained by the different intensity with
which the motion proceeds in the different organs. The unity of consciousness is
to be explained by the movements in the brain acquiring a direction that runs
back into itself; by making such circular movements possible, the brain is the
organ of consciousness. Thus Czolbe would conceive both sensation and
self-consciousness, plainly and simply, as movements in space. But --- clear and
consistent thinker that he is --- he becomes aware that if he identifies
sensation and self-consciousness with motion without more ado, then one must
also be able to reverse the proposition and say that wherever motion of a
certain intensity and a certain direction is present, there consciousness is
present too, so that the consequence is a universal ensoulment of nature. Thus
the consistent carrying through of materialism led beyond
% [p. 506]
it.\footnote{The consequence Czolbe draws from materialism was already drawn by
Henry More, who (\textit{De anima.} Roterodami 1677, p.~64) sets up the
following fundamental proposition: ``If sensation can be ascribed to matter,
then the motion of one material part, or its reaction upon another material
part, or at any rate the continuation of that motion, is one and the same as
sensation and apprehension. This proposition~... is confirmed by the voice of my
chief opponent, Hobbes.'' He then refers to Hobbes's doctrine of sensation as a
motion. Hobbes, however, denied that consequence, saying expressly that the
proposition that sensation is motion cannot be reversed. (\textit{De corpore}
XXV, 5).}

Later --- in the work \textit{Die Grenzen und der Ursprung der menschlichen
Erkenntniss} (1865), and in the exceedingly interesting paper \textit{Die
Mathematik als Ideal für alle andere Erkenntniss und das Verhältniss der
empirischen Wissenschaften zur Philosophie} (Zeitschrift für exakte Philosophie,
1866) --- Czolbe saw that it was impossible to explain existence from a single
principle, whether one finds that principle with Büchner in matter, or with the
speculative philosophers in spirit, or with the theologians in God. We reach an
explanation only by proceeding from several elements, calling that an element
which we are unable to analyse. Such elements, not reducible to one another, are
the material atoms, the organic forces, and the mental elements (which last make
up in sum the world-soul). Between these three kinds of element a harmonious
co-operation takes place, by which a purposive connection in nature is realized.
Not in its origin, but in its direction and tendency, does existence establish
itself as a unity. --- Thus it appeared that the problem of existence was more
intricate than materialism supposed. In his later investigations, which are only
partly published, Czolbe approached Spinoza's fundamental thoughts, of which he
wished to give an empirical execution. --- On the different stages of
development of this energetic thinker's philosophizing, reference must here be
made to \textsc{Vaihinger}'s paper \textit{Die drei Phasen des Czolbeschen
Naturalismus} (Philosophische Monatshefte XII).

What lies at the ground in Czolbe from first to last is the demand for clarity
and intuitability in all fundamental concepts. He wanted everything mystical and
supersensible away, and sought so far as possible to think all ideas with
geometrical clarity. It was more with respect to intuitability than with respect
to stringency of proof that he proclaimed mathematics the ideal of science. In
his youth he was enthusiastic for Hölderlin's poetry, and had set himself the
task of maintaining Greek clarity and joy in life against all mysticism and
dualism. Intuitability in the domain of thought was for him closely akin to joy
in the natural world. Just as it may require labour to carry through the
principle of intuitability in thinking, so it may cost resignation to carry
through joy in actual life in practice; but in the theoretical domain as in the
practical, Czolbe would combat what he calls ``der Blödsinn der Transcendenz''.
He admitted already in his materialist period that the innermost motive of his
philosophy was the subjective need for intuitability and for holding fast to the
actual; materialism was to him a postulate --- which was later replaced by other
postulates. ---

\textsc{Ernst Haeckel} (b. 1834, from 1865 professor of zoology in Jena) is
often reckoned to materialism, but has himself determined his standpoint as a
monism which leads beyond the opposition between spiritualism and materialism,
and which takes as its ground the great pantheistic thought of the unity of
nature. For monism, he says, there exists neither spirit nor matter in the usual
sense, but only one thing which is both at once. It was in his \textit{Generelle
Morphologie} (1862--1866) that he first uttered his general philosophical views.
The mental is for him --- as for Czolbe --- an original element in existence,
though it exists in exceedingly different degrees, from the atom-soul and the
cell-soul to the souls of the highest organisms. This theory of ensoulment
becomes questionable in Haeckel through his often letting the soul's action
serve as an explanation of organic movements, instead of seeking a purely
natural-scientific explanation. On the other hand he sometimes declares mental
phenomena to be mere effects of the most complicated and most unstable carbon
compounds, the albumen molecules in nerve tissue. He believes himself to be
beyond spiritualism and materialism, and yet one finds in him both propositions
that sound spiritualistic and propositions that sound materialistic. He has not
arrived at his principles through a sharp thinking-through of the laws of the
conservation of matter and of force. It is Darwin's doctrine that set his
thinking in motion, and he has sought a highest principle that could be united
with the great unity in nature to which the new doctrine pointed. Monism leads,
according to Haeckel (Gen.\ Morphol.\ II, p.~445--451), to the most exalted idea
of God. In the great, all-embracing law of causality the Godhead reveals itself
as embracing the whole of nature and operating in every natural phenomenon. The
usual conception of God, by contrast, is belief in two gods --- amphitheism, not
monotheism --- because it sets the natural causes beside the Godhead.

Haeckel was one of the first German investigators of nature to attach himself to
Darwin, and he fought with great ardour for his theory, to which he ascribed a
certainty, and gave an extension, that could not always be approved by the
theory's critical and level-headed founder. He was untiring in constructing
genealogical trees for the species now living, and he had no scruple about
assuming a constant arising of organic matter. He has no clear eye for the
limits of hypotheses and for the necessity of verification; only thereby is it
explained that he can without more ado set Darwin's hypothesis alongside
Newton's. He accordingly also reproached Darwin with laying too much weight upon
the objections to the theory. The great investigator shook his head at his young
adherent's zeal: ``Your boldness sometimes makes me tremble'', he wrote to him
(19 November 1868). --- In Haeckel's most recent ``philosophical'' work
(\textit{Die Welträthsel.} Gemeinverständliche Studien über monistische
Philosophie. Bonn 1899) the qualities mentioned appear in a still more glaring
form than in his earlier utterances. ---

The doctrine of the conservation of matter and of force, and the doctrine of the
development of species through natural selection, were --- as already remarked
--- the natural-scientific results which in the period after the middle of the
century had above all to set reflection in motion. Having now considered the
attempts made by investigators of nature to found a philosophy upon these
results, we shall go on to examine the way in which professional philosophers
took up their position with respect to them.

\section*{3.\quad Idealist Constructions on a Realist Basis}
\label{sec:idealistiske}

\subsection*{a.\quad Rudolph Hermann Lotze}

% [p. 518]
Idealist philosophy has in Lotze its most important representative in the last
half of the century. In his personality and in his development he stands as one
who has taken up into himself, in an interesting way, both the ideal motives
upon which the philosophy of Romanticism built and that strict carrying through
of the mechanical view of nature which towards the middle of the century was
working its way forward in science. Lotze is a master of conceptual development
--- of discussing a thought so that all its nuances come out, of taking up a
problem again and again and treating it from different points of view. His ideal
was in the last instance the same as the one that floated before the philosophy
of Romanticism: to derive all development and all connection in the world from
an eternal Idea, which should contain the ultimate ground of all that happens,
as well as of the value of what happens. But he saw clearly that such a
derivation exceeds what human thought can accomplish, and that when the thinkers
of Romanticism nevertheless believed they had given one, it was because in them
the poetic-religious tendency involuntarily merged with the philosophical. Lotze
therefore separates the poetic-religious element from the speculative. His own
fine poetic sense let him understand that their Romantic confusion was not a
happy one. And this poetic sense is one with a sense for individual nuances and
relations, which speculative philosophy all too often dissipated by its
abstractions. Through this sense Lotze's speculative interest hangs together
with his realist tendency. What matters for him is first and foremost to
apprehend the phenomena in their concrete nature and in their determinate,
law-governed connection; it is then the task of philosophical thinking to find
out the presuppositions upon which that real connection rests. The poetic, the
natural-scientific and the philosophical element in Lotze thus hang closely
together, and he was equipped as few have been to treat the task he had set
himself, and which the intellectual conditions of the age suggested: to attempt
a reconstruction of an idealist philosophy upon a realist basis. It was his
conviction that the philosophy of Romanticism had overlooked the real conditions
and the mechanical connection of nature, without which even the most significant
ideas are helpless idols; just as it was his conviction that materialism made
\emph{that} into the be-all and end-all which was only a --- quite certainly
necessary --- form and frame for the valuable content of existence. His
philosophy is essentially an analysis of the very concept of natural mechanism,
designed to show that this concept leads with necessity to the assumption of an
ideal principle of existence, and that it at any rate does not exclude the
assumption that such a principle is the eternal source of all that is good and
valuable in the world.

In Lotze's development the different motives and interests that animated him
appear clearly. He was born in Bautzen in Saxony on 21 May 1817, in the same
region from which Lessing and Fichte come. At the university of Leipzig he
studied philosophy, medicine and physics. He was here initiated into the two
circles of thought which it was later to be his task to unite, after having
worked through them separately. It was the aesthetician and philosopher of
religion Chr. Hermann Weisse who was his teacher in philosophy, and Lotze later
said that he owed to Weisse not merely many particular influences, but that he
was in particular introduced by him into a circle of thoughts which he has since
felt no need to give up. Weisse was the most eminent representative of
philosophical theism. Through him Lotze's religious-philosophical conception
stands in historical connection with Schelling's --- and through that with old
Jacob Böhme's. Lotze emphasizes, in a retrospect upon his development
(\textit{Streitschriften.} 1857, p.~5\,ff.), that he appropriated the ideas of
speculative philosophy not as a closed doctrinal edifice but as a peculiar form
of intellectual culture. The presupposition that the ultimate ground of things
could be thought only as spiritual was early laid down in him and never given
up. Medicine and physics he studied under E. H. Weber, Volkmann and G. Th.
Fechner. Here he came to know the natural-scientific method and conception at
first hand. In one and the same year he took the doctorate in philosophy and in
medicine, and now came forward as lecturer in both subjects --- also after a
professorship in philosophy had been conferred upon him in 1842. After working
some years in Leipzig he became professor in Göttingen, as Herbart's successor,
and here fall his most important works. The year before his death he received a
call to Berlin, but there soon succumbed to an illness from which he had long
suffered (1881). His life was quiet, devoted to study, thought and academic
teaching. With a rare union of many-sidedness and thoroughness he knew how to
orient himself in highly different fields --- something to which not only his
medical and philosophical works but also a whole series of smaller papers and
reviews (published after his death in four volumes under the title \textit{Kleine
Schriften}) bear witness. From his more rigorous scientific labours he rested in
occupation with art and literature: thus he occasionally translated ``Antigone''
into Latin verse.

As a medical writer Lotze set himself, as already touched upon, the task of
maintaining physiology's character as a mechanical natural science. He would not
have the peculiarity of the organic phenomena explained by appeal to a
mysterious vital force, but by exhibiting the determinate, law-governed manner
in which the general forces of nature operate within the organism. As everywhere
in nature, so here too we must find the explanation in the interaction of real
elements. Not by being withdrawn from the mechanical connection of nature, but
by the peculiar manner in which connected series of effects arise, does organic
life distinguish itself from the inorganic world. In speculative philosophy of
nature the idea of the organism had been made the type for the whole conception
of the world; that this idea was corrected had therefore to be of general
philosophical significance, and one understands that not only Lotze's medical
but also his philosophical interest was satisfied by his labours in this
direction (\textit{Allgemeine Pathologie und Therapie als mechanische
Naturwissenschaften.} 1842. --- \textit{Allgemeine Physiologie des körperlichen
Lebens.} 1851). These works even gave some materialists occasion to see a
fellow-thinker in Lotze, notwithstanding that he stated definitely that
mechanism was for him only a part of his conception of nature, not the whole. It
is characteristic of Lotze that he thus divides the problems. His general
philosophical views he had developed in his \textit{Metaphysik} (1841). But he
by no means felt himself finished when he began his teaching activity, and
especially as regards the concluding ideas he constantly referred to future
works. In \textit{Medicinische Psychologie oder Physiologie der Seele} (1852) he
expressed himself in detail upon the relation between the spiritual and the
material, and at the same time gave psychological investigations of lasting
interest. A long-cherished plan he carried out in the following period
(1856--1864) with the publication of his \textit{Mikrokosmus} (3 vols.). This
work, conceived as a counterpart to Humboldt's ``Kosmos'' and Herder's
``Ideen'', gives a psychology in close connection with physiology and cultural
history, and ends with a development of Lotze's cosmological and
religious-philosophical ideas. For all his interests, for all the different
roads his thought and feeling had travelled, he had here found expression. The
work is laid out as a popular presentation and has won a considerable
circulation. After publishing a spirited \textit{Geschichte der Aesthetik in
Deutschland} (1868), Lotze proceeded to give a concluding systematic
presentation of his philosophy. But only two parts of it did he succeed in
finishing (\textit{Drei Bücher der Logik.} 1874. --- \textit{Drei Bücher der
Metaphysik.} 1879). The third volume, which was to comprise aesthetics, ethics
and philosophy of religion, was not worked out. His philosophy is thus
unfinished; it was not granted him to set the crown upon the work and show how
the many threads of thought he had spun could be woven together into unity. ---
Short presentations of Lotze's teaching in the different philosophical fields
are to be had in the dictated notes (\textit{Grundzüge}) from his lectures,
published after his death. Among these must especially be singled out
\textit{Grundzüge der Logik und Encyclopädie der Philosophie}, which (in its
second part) gives a short and clear survey of his whole system. Here we give a
characterization of the three most important points in his philosophy.

\subsubsection*{α.\quad The Mechanical View of Nature}

Lotze's thinking has two points of departure: a living feeling for the value of
spiritual life, for the fact that the highest thing for us is bound up with
spiritual development and its ideals --- but at the same time a firm conviction
that a system of mechanical causes and laws is necessary to the realization even
of the highest ideals. He can, as he says somewhere, believe in ideas that
operate, but not in ideas that bewitch. The two points of departure do not,
however, stand wholly isolated in him. He seeks, by proceeding from the one, to
show the possibility of acknowledging the other. The most characteristic and the
most scientific feature of his philosophy is his attempt to establish, through an
analysis of the very concept of mechanism, that its ultimate presupposition is
the concept of a principle which --- as soon as we conceive it in a way
intelligible to us --- shows itself as the bearer of, and the source of, the
highest ideas. Lotze means to take the bull by the horns: by drawing out
realism's consequences he will found idealism. Only investigations conducted in
realism's spirit can, in his view, bring us nearer to the goal idealism sets
itself: the knowledge of the world as the expression of a valuable idea.
Romanticism's philosophy had believed it could derive the forms of reality from
the highest Idea. That has shown itself to be an insoluble task. We can, by
contrast, attempt by way of inference to go back from what is given to its
presuppositions. Not a \emph{deduction} but a \emph{reduction} is possible.
Thinking must always apply its forms to a given material. Both general culture
and the several sciences here operate with a mass of concepts whose origin,
meaning and validity they do not investigate more closely. Such concepts are
cause and effect, matter and force, end and means, freedom and necessity, matter
and spirit. It is philosophy's task to bring unity and connection into the circle
of representations by making these concepts --- presupposed in practical life and
in the several sciences --- the object of a peculiar investigation, and by
determining the limits of their validity (\textit{Grundzüge der Logik und
Encyclopädie der Philosophie}, §~88).

Of these concepts the most important is the one we presuppose in every
investigation of reality --- every time we make an experiment, every time we seek
an explanation: the concept of a pervasive, all-embracing causal connection. This
concept is not grounded upon experience, but is presupposed in all experience.
But since the understanding of nature so far won rests upon its reality, it can
itself be described as the expression of a fact; and that fact --- that the
single element in our experience is bound by a law-governed connection to other
elements, the fact of mechanical connection --- is what philosophy has to think
through and to draw all the consequences of. Deduce it, philosophy cannot; but it
may perhaps discover what lies in it. A \emph{multiplicity of real elements in
mutual interaction} is the foundation upon which the mechanical view of nature
builds; and Lotze was, as we have seen, so convinced of that foundation's
inevitability that he himself worked with zeal to gain it recognition within
physiology, where hitherto vitalism had reigned with its appeal to a single
forming and directing ``vital force''.

But because mechanical connection is an inevitable feature of our conception of
the world, it need not be the only, all-determining feature. On the contrary:
even if its validity is unlimited, it may perfectly well have a subordinate
significance in the world as a whole. (\textit{Die unbeschränkte Gültigkeit des
Mechanismus, aber zugleich seine durchaus untergeordnete Bedeutung im Ganzen der
Welt!} \textit{Drei Bücher der Metaphysik}, p.~462.) And a closer analysis of the
concept of mechanism leads us to see that this is precisely how the matter
stands.

The mechanical view of nature stops --- when it is made into a concluding
world-view --- at a multiplicity of elements (atoms) in mutual interaction. It
proclaims a pluralism. But how do the elements stand to the connection in which
they stand? Can they have independent subsistence apart from that connection, so
that it is a relation indifferent to their essence --- or are they not determined
right through by the combination in which they stand with the world as a whole?
Interaction and connection cannot, after all, float freely above or between the
elements, but presuppose an inner unity of them. For if I assume that the two
elements \textit{A} and \textit{B} are wholly independent, interaction between
them becomes unintelligible. An effect cannot be transferred as a fully finished
state from \textit{A} to \textit{B}. What goes on in \textit{A} can acquire
significance for what goes on in \textit{B} only if \textit{A} and \textit{B} are
not in reality independent, absolutely separated beings, but if their states are
states in one and the same, more comprehensive being. A multiplicity of
independent beings would make mechanical interaction unintelligible;
intelligibility comes only with the assumption of an infinite, all-embracing
being, of which the single elements are moments, or points of action. The
concept of a force's or an influence's ``passing over'' from one independent
element to another is untenable. Not a \emph{transient} cause (\textit{causa
transiens}), only an \emph{immanent} cause (\textit{causa immanens}) is
intelligible. The states of one and the same being can stand to one another as
ground and consequent; the states of two mutually independent beings cannot. ---
Among the numerous presentations Lotze has given of this train of thought,
\textit{Logik und Encyclopädie} §§~99--100 and \textit{Drei Bücher der
Metaphysik} §§~50--81 must be singled out.

By analysis of the concepts of the causal relation and of interaction --- the
fundamental concepts of the mechanical view of nature --- Lotze has thus arrived
at the idea of a primal substance, an all-embracing principle. His thinking here
moves along the same paths as those which in their day led Spinoza to his concept
of substance, and which Kant walked especially in his youthful writings, without
however abandoning them in his later works. He is reminiscent of Leibniz too at
several points, without there seeming here, any more than in Spinoza's case, to
have been any influence through study. When he heard that he was being placed
alongside Spinoza and Leibniz, he said (see \textsc{Falckenberg}:
\textit{Hermann Lotze}. Stuttgart 1901, p.~85) that he did not feel himself to
stand in any inner relation to them, though he admitted that there was occasion
enough to institute a comparison.

The concept which Romanticism's philosophy would lay at the ground, and from
which it would derive everything, becomes for Lotze the last postulate --- or, as
he also expresses it, the last fact for our thinking. A closer, intuitable
development of the universal principle presupposed in the simplest relation of
interaction is impossible. We stand here at a \emph{limiting concept} which we
can neither do without nor develop. (\textit{Drei Bücher der Metaphysik},
§§~73 and 246.) It is nevertheless this concept that makes it possible for Lotze
to hold fast to what was for him the essential thing in the idealist train of
thought into which his teacher Weisse had introduced him. By it he found a
connection between the opposed currents of the world of thought. Neither
materialism's absolute atoms, nor Leibniz's monads, nor Herbart's reals could for
him mark the conclusion of thinking: pluralism must necessarily give place to
\emph{monism}. As regards the atoms in particular, Lotze maintains that natural
science's interest leads only to assuming elements which are in fact indissoluble
for our experience, but that the assumption of a multiplicity of extended
elements --- even if these are thought of as vanishingly small --- cannot be any
conclusion for thought. Either one must give up the atom's unity or its
extension: within an extended atom every effect will take time and pass from part
to part --- and these parts will then be more fundamental unities than the whole
``atom''. In the concluding concept of the atom we must abstract from all
extension and think of the atoms as centres of force which --- according to the
analysis of the concept of mechanism given above --- are each of them points of
departure for the primal substance's operation. (\textit{Drei Bücher der
Metaphysik}, §§~190--191 and 245; cf.\ the interesting review of Fechner's
``Atomenlehre'' in \textit{Kleine Schriften} III, p.~215--238.)

Lotze's assertion of the mechanical view of nature, and the consequences he draws
from it, make up the most significant part of his philosophy --- notwithstanding
that it is as a rule not brought forward when Lotze is praised. It is the
spiritualistic and theologizing tendency in him that is most frequently
emphasized, by admirers and opponents alike. And yet he has earned his greatest
merit as a thinker by his analysis of the fundamental concept of the scientific
view of nature. It is a defect in him that he does not undertake an
epistemological investigation of the limiting concept he arrives at, but lets it
stand in the form in which Romanticism's philosophy and the earlier dogmatism had
conceived it. Lotze had in general no understanding of the importance of
epistemology.

\subsubsection*{β.\quad Metaphysical Idealism}

In Lotze's train of thought, so far as we have presented it hitherto, there is
still a certain indeterminacy. It might still be asked whether we then have no
way at all of forming a closer representation of the elements and of the primal
substance. In now bringing forward Lotze's answer to this question, it must first
of all be emphasized that he is entirely clear that any answer to it rests upon
analogies, and is determined by motives other than the purely theoretical. The
law of analogy, and the need to find our own spiritual being again in existence
--- which have shown themselves to lie at the ground in all metaphysical
idealists, and especially in Leibniz, Berkeley, Schelling, Schopenhauer and
Beneke --- appear with full consciousness in Lotze.

Lotze is an atomist --- but he does not conceive the atoms as material, because
extension, as well as all other sensible qualities, must be explained by the
interaction of atoms, which therefore cannot themselves have that property. Like
life and all qualities of experience, the sensible fact of extension is due to an
interplay of force-points, which are in turn to be conceived as initial points of
the infinite primal being's inner operation. Lotze does not regard it as
impossible that wholly new beginnings may occur in the course of the world: this
does not conflict with universal lawfulness, for the law always expresses only
the order in which the states arise, not an outer fate. The number of the
elements need not be constant, but may be thought to rise or fall according to
the demands of the highest world-purpose. (\textit{Drei Bücher d.~M.} §§~33 and
235.) Nor need the single elements be wholly homogeneous with one another. A
certain agreement or commensurability must indeed be present if one world-order
is to embrace them all; but a complete likeness is not necessary. A law of nature
may very well bind together qualitatively different elements for which we are
unable to find any common denominator. It is in the end not a \emph{logical} but
an \emph{aesthetic} necessity that makes existence intelligible to us
(\textit{Drei Bücher d.~M.} §~59): not a formal consistency in the primal being's
operation, but the richness and fullness with which that operation proceeds, is
decisive for the constitution of the single elements. At this point Lotze plainly
enough passes over from purely theoretical motives to motives of feeling ---
without either the one or the other enabling him to determine more definitely how
great \emph{differences of quality} are compatible with the
% [p. 518]
\emph{commensurability} which the fact of interaction
presupposes.\footnote{On the indeterminacy in Lotze at this point see Max
Wentscher: \textit{Lotzes Gottesbegriff und dessen metaphysische Begründung.}
Halle 1893, p.~19--29.}

If we would now form a representation at all of the inner nature of the elements,
we must conceive them on the analogy of our own spiritual being. The mechanical
view of nature really instructs us only about the mutual relation between the
elements, not about their inner nature. It concerns the outer circumstances,
moves the whole time around the outside of things --- is a \textit{cognitio circa
rem}. It is perhaps even inclined --- like the popular conception of the world
--- to hold that it is wholly indifferent to things whether we apprehend them or
not. In opposition to this Lotze maintains (already in his youthful work
\textit{Metaphysik}, 1841, p.~313) that subjectivity is just as much a part of
reality as the outer objects are: \textit{Nicht nur was zwischen den Wesen,
sondern auch was in ihnen sich zuträgt, ist ein wahres und wirkliches Geschehen}.
In his later writings --- first in \textit{Medicinische Psychologie} --- he goes a
step further, maintaining that our own subjective being is the only case in which
we know a thing's interior and have a \textit{cognitio rei}, not merely a
\textit{cognitio circa rem}: the only way of forming a representation of the
interior of other things is then to conceive them on the analogy of ourselves ---
that is, as feeling beings (if not as representing beings, since feeling is a
more original expression of consciousness than representation). The use of this
analogy is the only means of thinking things as real beings subsisting for
themselves, which are more than mere images for us. We thereby merely make use of
the general method that leads the unknown back to the known. (\textit{Drei Bücher
d.~M.} §§~96--98.) And the ``aesthetic necessity'' demands this too, since it
cannot be reconciled with an existence in which a great part --- perhaps the
greatest --- should be merely a dark foundation for a mental life extended to
only a few regions (see already \textit{Allg.\ Physiologie} §~129). The elements
of existence must therefore be assumed to be ensouled in manifold degrees.

As it stands with the elements of existence, so it stands also with the
all-embracing world-ground that makes their interaction intelligible. If we
follow the law of analogy and attempt to lead the unknown back to the known, then
according to Lotze we must conceive the world-ground as a spiritual being. When
inner states are to be bound into unity and connection, our own inner, spiritual
life is the only example we know in which the possibility of such a preservation
of unity through the different states is realized --- at least approximately
realized. Following his teacher C. H. Weisse, Lotze conceives the world-ground as
an absolute personality, and he defends this transference of the concept of
personality to the absolute being on the ground that only an absolute being can
be a personality, because only it has inner independence and originality, while
the concept of personality finds only imperfect application to finite beings,
which are dependent upon outer conditions. Lotze admits, to be sure, that to a
personal life there belong resistance to be overcome, and a capacity to suffer
and to receive as well as to act. But when one then asks how an absolute being,
subject to no limits, can suffer, Lotze answers (\textit{Grundzüge der
Religionsphilosophie} §~34, cf.\ §~52) that what sets the Godhead's feeling in
motion must be ``inner products of its own creative imagination''! --- It is
surely a great question, however, whether such a self-created resistance can mean
anything serious --- especially as it can be annihilated at whatever moment one
pleases. The personal life we know has in any case precisely not self-created and
not so easily removable limits to struggle with, and the analogy upon which Lotze
relies therefore seems to fail at just the decisive point. And to this it must be
added that Lotze --- on the most probable interpretation of his obscure and
% [p. 520]
shifting utterances\footnote{Reinhold Geijer, in his paper \textit{Hermann Lotzes
tankar om tid och timlighet i kritisk belysning} (Lund 1886), drew attention to
the shifting standpoints Lotze seems to have taken on the question of time's
absolute validity. In my paper \textit{Lotze og den svenske Filosofi}
(Letterstedtske Tidsskrift 1888) I associated myself with Geijer's
interpretation, though I had to conceive the problem itself otherwise than my
Swedish colleague. Yet the matter seems to stand otherwise than Geijer, and I
with him, had reason to assume. According to what Richard Falckenberg has shown
in his paper \textit{Die Entwickelung der Lotzeschen Zeitlehre} (Zeitschrift für
Philosophie und philosophische Kritik, vol.~105), the \textit{Grundzüge der
Metaphysik} published after Lotze's death date from as early as 1865, and with
that an essential support falls away for the assumption that Lotze first denied
the reality of time towards the close of his life. When in earlier writings he
speaks without qualification of the time-relation as absolutely valid, this is
according to Falckenberg a legitimate use of popular representations where the
context did not require a more stringent usage. Lotze's definitive conception was
the one presented in \textit{Drei Bücher der Metaphysik} (1879) and in
\textit{Grundzüge der Religionsfilosofi} (from 1878--79), according to which
succession holds for finite beings, while the Godhead is raised above every
difference of time. It departs from his earliest presentation (in
\textit{Metaphysik}, 1841) in asserting this distinction, which admittedly does
not make the whole question clearer --- though the \emph{psychological}
distinction between succession and abstract time-form may well be defensible.
--- As little as my Swedish colleague am I able to connect any meaning with the
expression ``ein zeitloses Geschehen und Wirken'' (Grundzüge der Metaphysik
§~58). My German colleague, by contrast, finds a significant thought in it.} ---
differs from Weisse in that he does not hold the form of time applicable to the
absolute being. A personal being that does not develop in time --- a timeless
life, and a timeless suffering and acting: these are concepts that make far too
great demands upon our capacity for analogy! --- It need only further be remarked
that Lotze --- as regards both the elements and the world-ground --- presupposes
that the opposition between spirit and matter is contradictory, so that no
further possibilities could be thought than the two: only on that presupposition
is there compelling necessity to conceive the elements and the world-ground as
spiritual beings, when one cannot or will not conceive them as material. But
there is no logical, only a factual necessity for choosing between spirit and
matter in determining the concluding concepts. We in fact know only these two
forms of existence; but there would be nothing \emph{unthinkable} in one or more
other forms existing. Therefore metaphysical idealism cannot be
proved.\footnote{Compare my \textit{Psykologi}\textsuperscript{4}, p.~80;
\textit{Psykologiske Undersøgelser}, p.~91\,f.; \textit{Filosofiske Problemer}
(Universitetsprogram 1902), p.~62--65. --- It is interesting that the founder of
metaphysical idealism, Leibniz, in his earlier period, like Spinoza, laid great
weight upon the duality spirit--matter being purely factual, not logically
necessary. \textit{Philos.\ Schr.}, ed.\ Gerhardt I, p.~237, 242, 268. Compare
Ludw.\ Stein: \textit{Leibniz und Spinoza}, p.~93\,f.}

Lotze is now also fully conscious that the conclusion of his whole world-view
rests upon motives of feeling, not upon strict thinking. The necessity of going
back from the multiplicity of the elements to the unity of the world-ground he
asserts, indeed, on purely theoretical grounds: ``\textit{Gar kein Weltlauf,
weder ein harmonischer noch ein unharmonischer, ist mir ohne die Voraussetzung
jener Einheit begreiflich, die alles gegenseitige Wirken des Einzelnen erst
möglich macht; die Störungen der Dinge durch einander bezeugen die beständige
Gegenwart dieses Einen eben so eindringlich, wie das Zusammenstimmen der Kräfte
zum Zweck}'' (\textit{Drei Bücher d.~M.} §~233). But when he leads everything
back to a single principle, he knows very well that we cannot construct the
world-connection and the elements out of it: we always lack the second premiss!
(\textit{Drei Bücher d.~M.} §~93.) Only in practical conviction, then, can the
thought be held fast that all beings and all events have their ultimate ground in
the highest purpose maintained by the world-ground. What that purpose is we do
not know --- any more than we know how the conflict of forces in nature and the
reality of evil in the moral world are compatible with the validity of the
world-plan. (\textit{Mikrokosmus}, Buch IX, Kap.~5. --- \textit{Drei Bücher
d.~M.} §~233.) But in an ethical idea Lotze found, at the beginning as at the
close of his career, the final conclusion for his thinking: metaphysics goes back
for him to ethics; in what \emph{ought} to be lies the ultimate explanation of
what \emph{is}. ---

If by pantheism one understands the doctrine which maintains an inner relation
between God and the world --- whether the Godhead is conceived as a personal
being or not --- then Lotze's religious-philosophical standpoint is to be
described as ethical pantheism. He himself felt closely akin to Fichte; the
kinship with Spinoza he would not properly acknowledge. In the use of theological
words and turns of phrase he went --- especially in \textit{Mikrokosmus} and in
the \textit{Grundzüge} of philosophy of religion and of moral philosophy
published after his death --- further than was fitting and warranted. By his
usage (for it is nothing more) it came to look as though he stood nearer to the
popular religious circle of ideas than in fact he did. He does, however,
occasionally emphasize (especially \textit{Mikrokosmus}, Buch VIII, Kap.~4) the
difference between free religiosity and the positive religions.

\subsubsection*{γ.\quad Spiritualistic Psychology}

The elements whose interaction produces the phenomena that experience presents
need not --- as we have seen --- according to Lotze be absolutely homogeneous.
There need be only such homogeneity as the interaction itself presupposes. By
this very indeterminate rule Lotze makes it possible for himself to hold fast to
the popular assumption of an interaction of soul and body. The elements which the
concept of mechanism presupposes are for him partly psychical, partly physical.
He disputes the assumption of a continuous series of physical phenomena. At
certain points the physical motion is broken off, to be ``absorbed'', to pass
over into psychical states. It is indeed doubtful whether Lotze, when he first
put forward this claim (in the year 1851; \textit{Allg.\ Physiol.} §~424), had
any living conception of the significance of the newly discovered law of the
conservation of energy --- though neither did he later find in it any reason to
depart from the standpoint he had once seized. There were for him compelling
grounds which made the assumption of a peculiar soul-substance necessary. Only
two possibilities presented themselves to him: either to derive the psychical
phenomena from a soul as a principle peculiar to them, or to explain them by the
co-operation of the physical forces. (See \textit{Selbstanzeige der medicinischen
Psychologie}, Kleine Schriften III, p.~4.) Since the latter is impossible ---
inasmuch as the interaction of physical forces cannot explain the unity that
stamps even the simplest expression of mental life --- only the first way out
remains for Lotze. A great part of Lotze's psychological investigations is
therefore devoted to finding out the manner in which the psychical elements stand
in interaction with the physical ones. His psychology here recalls Descartes's.
There are, on his view, certain regions in which the physical elements act upon
the soul; thereafter the mental processes may for a time run their course
exclusively according to their own laws, and they may then in turn release a
mechanical force for the production of new physical changes. (Kleine Schriften
III, p.~7.) Notwithstanding that the present condition of nerve physiology and of
psychology would seem to make a separation of these different stages hopeless,
Lotze nevertheless proceeds confidently (especially in \textit{Medicin.\
Psychol.} and \textit{Mikrokosm.}) to a closer investigation. The significance
which the material processes have for mental life cannot --- according to the
conception Lotze has urged of the nature of interaction --- consist in finished
sensations or representations being conveyed from them to the soul; they can only
give signals, which the soul must translate into its own language, just as,
conversely, the soul's inner states are not as such transferred to the material
organs, but only become the occasion of these entering into function. The
material organs work in the service of the higher spiritual activities by
delivering and preparing the material upon which the soul is to exercise its
powers, and in that their whole significance consists. The higher spiritual
activities themselves --- memory, thinking, aesthetic and moral feeling --- take
place, once the material is given, in the soul's interior, and it is wholly
unnecessary to assume peculiar material processes corresponding to them. Lotze
would in any case sooner interpret such material processes as effects of the
mental states --- brought about by these vibrations propagating themselves from
soul to brain --- than as their causes. So especially when we see a state of
feeling persist long after the original occasion has fallen away.

Upon Lotze's special psychological theories we cannot here enter. Particular
mention must be made of his brilliant doctrine of local signs, by which the
development of the apprehension of space becomes more
% [p. 523]
intelligible,\footnote{See on this Reinhold Geijer: \textit{Hermann Lotzes lära
om rummet} (Nyt svensk tidskrift 1880) and \textit{Darstellung und Kritik der
Lotzeschen Lehre von den Localzeichen} (Philosophische Monatshefte 1885). Cf.\
my paper \textit{Lotze og den svenske Filosofi} (Letterstedtske Tidsskrift 1890)
and my \textit{Psykologi} V~C,~7~c.} and of his assertion of feeling's importance
as a fundamental element in mental life, over against Hegelian and Herbartian
psychology. On a mass of particular points Lotze throws new light by his fine,
spirited conception and his often happy manner of expression. It conflicts with
good psychological method, on the other hand, when Lotze, instead of seeking the
laws that connect the mental phenomena with one another, conceives the relation
in the arising of the different mental phenomena in such a way that the
representations react upon ``the soul itself'' and release the feelings out of
it, which must in turn act upon ``the soul itself'' in order to release from it
the expressions of will. Likewise when he lets a peculiar faculty of thinking
appear in opposition to the association of representations. Here his
spiritualistic metaphysics has forced upon psychology an all too artificial and
scholastic mode of explanation. It takes its revenge here that Lotze does not
acknowledge psychology as an independent science of experience, but regards it as
applied metaphysics. (\textit{Logik und Encyclopädie} §~93.)

When Lotze's psychology recalls Descartes's on essential points, it must
nevertheless be held fast that this likeness is present only so long as Lotze
holds to the popular usage which he employs, in his psychology no less than in
his philosophy of religion, to a great extent. The soul a thing for itself and
the body a thing (or group of things) for itself --- from many of Lotze's
utterances that might seem to be his doctrine, as it was Descartes's and still is
that of popular metaphysics. But it is only a provisional conception for Lotze.
From Descartes he differs first and foremost by his doctrine of the subjectivity
of extension and by his analysis of the concept of matter, which leads him to the
assumption that matter is only a phenomenal form of the interaction of psychical
elements. Thereby the dualistic conception of the relation between soul and body
is abolished. The qualitative difference between the interacting elements is
thereby considerably reduced: instead of an interaction of psychical and physical
elements we get a mutual interaction of psychical elements. And this reduction
Lotze makes --- as must be well noted --- not because he finds any difficulty in
popular or in Cartesian metaphysics. The problem of the relation between soul and
body really does not exist for him at all; the interaction between them presents
to him no greater riddle than the relation of interaction that takes place
% [p. 524]
between two bodies striking one another.\footnote{\textit{Mikrokosmus} Buch III,
Kap.~1: ``\textit{Warum sollte nicht ein Atom des Nervensystems ebenso auf die
Seele oder sie auf jenes stossen und drücken können, da doch jeder gemeine Stoss
und Druck sich für die nähere Betrachtung nicht als ein Mittel zur Wirkung,
sondern nur als die anschauliche Form eines viel zarteren Ereignissen zwischen
den Elementen ausweist?}''} It is general philosophical considerations that guide
him here --- considerations that would retain their validity even if he had,
instead of laying Cartesian dualism at the ground of his psychology, built upon
Spinoza's identity hypothesis. Metaphysical idealism is independent of every more
special psychological hypothesis. Since Lotze began as a Cartesian, those
considerations lead him to a conception that may be described as a
\emph{monistic spiritualism}, in that it regards soul and body as two
\emph{different substances} which are nevertheless of \emph{homogeneous}
constitution. Herbart's theory, and a theory that appeared within the Wolffian
school (see Note~1 to this volume), form predecessors of this peculiar theory.

On another point, and one more essential for his whole philosophy, Lotze also
separates himself from Descartes --- and here from Leibniz and Herbart as well.
As we have seen, he regards the single elements in existence as moments or
expressions (actions) of the primal substance, which is the one thing that is in
truth real. The soul can then, no more than the physical atom, be substance in
the stricter sense of the word. The independence of finite things is only
apparent. ``Our monistic conception'' --- says Lotze --- ``has \emph{wholly and
without reserve} laid the ordering of the world, and the existence and efficacy
of every single thing, in the hands of the one infinite being upon which alone
the possibility of all interactions rested.'' (\textit{Drei Bücher d.~M.} §~245,
cf.\ §§~72--73.) The words substance and being he therefore takes, when he uses
them of the soul, only in the sense of that which can act and suffer (ib.\
§~243) --- and he admits at last (ib.\ §~307) that it might perhaps be best to
refrain from using these words, which can easily mislead. Already in his paper
\textit{Seele und Seelenleben} (in Wagner's ``Physiologisches Handwörterbuch'',
1846; ``Kleine Schriften'' II, p.~198) Lotze said that every soul possesses only
the reality that belongs to it on account of its significance in the whole of the
world, and that its immortality depends not upon its nature but upon its place
within the ethical world-order. In ``Mikrokosmus'' and in his principal work
(\textit{Drei Bücher der Metaphysik}) Lotze returns to the same thought. Just as
it was not an interest in the soul's immortality that led him to his
spiritualistic psychology, so he sees clearly that this question cannot be
answered by scientific means at all. ``No other fundamental proposition stands
here at our disposal than the general conviction: all that has come to be will
subsist, if its subsistence belongs to the world-plan, and so long as it belongs
to it; all will perish whose reality had its warranted place only in a passing
phase of the world's course. That this fundamental proposition cannot be applied
in detail by human hands scarcely needs saying; we certainly do not know the
merits that could give one being a right to eternal subsistence, any more than
the defects that could deny another that right.'' (\textit{Drei Bücher d.~M.}
§~245.)

When one notices the agreement between Lotze's earlier (1846 and 1864) and later
(1879) utterances upon this point, one will scarcely be led to assume, with some
% [p. 526]
recent critics of Lotze,\footnote{Kresto Krestoff: \textit{Lotzes
metaphysischer Seelenbegriff.} Halle 1890, p.~46\,f.; 73. --- Max Wentscher:
\textit{Lotzes Gottesbegriff und dessen metaphysische Begründung.} Halle 1893,
p.~11\,ff. --- The opposition in which Lotze's conception stands from the very
first to Herbart's, despite all the points of contact it has with it, also
testifies that no change of principle has taken place in Lotze's thinking in the
course of the years. Compare on this Max Nath: \textit{Die Psychologie Hermann
Lotzes in ihrem Verhältniss zu Herbart.} Halle 1892. --- With Lotze's ethical
monism his optimism hangs together. Although he himself repeatedly emphasizes the
reality of physical and of moral evil as an obstacle to the carrying through of
his world-view, an otherwise well-disposed critic has nevertheless complained
that he lacks interest in ``the fundamental dogma of all religions, the
atonement'', because he does not take sufficient account of evil. G. Vorbrodt:
\textit{Principen der Ethik und Religionsphilosophie Lotzes.} Dessau und Leipzig
1891, p.~97.} that Lotze made essential alterations in his doctrine of the soul's
substantiality and thereby approached Spinozism more nearly than before. But it
would surely have been to the advantage of clarity if the excellent thinker ---
whose disposition assuredly permitted no unworthy accommodation --- had already
in his earlier presentations brought out his divergence from the current
assumptions as definitely as in the latest utterances.

\subsection*{b.\quad Gustav Theodor Fechner}

Fechner and Lotze stand as two Dioscuri in German philosophy in the last part of
the century. They are kindred spirits in their ideal disposition, their great
natural-scientific culture, their poetic sense and their need for unity in the
conception of the world. They pursue kindred ends, but in part by different
roads. A comparison between them is therefore interesting, especially as on
decisive points one must choose between them. The choice will --- since the
fundamental proposition, an idealist world-view upon a realist foundation, is the
same in both --- depend upon which of them lets the realist foundation come most
fully into its own. That is a measure Lotze has himself acknowledged, in the
remark that idealism can be maintained only by investigations conducted in
realism's spirit. Even if Lotze takes the first place as regards the
working-through of the general philosophical principles, Fechner may well have
the precedence as regards the consistent and resolute holding fast to the
natural-scientific foundation. This is all the more remarkable in that Fechner's
imagination has from the first a far greater influence upon his thinking than is
possible in his critical, deliberate companion. Fechner is --- like Kepler, of
whom he strongly reminds one --- an interesting example of how one may be led
through bold and fantastic speculations to positive and exact results, provided
only one holds fast to a definite fundamental thought and is able to extract it
from its mystical wrapping. Just as Kepler was led successively from mystical
speculations to the discovery of his famous laws, which satisfied his need to
find determinate mathematical relations expressed in the universe, so Fechner was
led through bold analogies to the conviction of a determinate quantitative
relation between the spiritual and the material, and, in working that thought out
more closely, came to found psychophysics, or experimental psychology.

Fechner was born in Lausitz in 1801 and in his youth pursued medical and physical
studies. In 1835 he became professor of physics in Leipzig and made himself known
by able work in his subject. The philosopher of religion C. H. Weisse belonged to
his closest friends and had great influence upon his religious conception --- as
upon Lotze's. It is nevertheless doubtful whether Fechner would have come to play
so great a part in the history of philosophy as he has, had he not, by studying
subjective phenomena of light and colour (in the winter of 1839--40), ruined his
eyes, so that after long suffering he had to give up his professorship of
physics. In philosophical thinking, in free fancy, in absorption in the inner
world of the mind, he sought compensation for the bright and colourful world in
which his weakened sight no longer allowed him to dwell as persistently as
before, but which he did not forget. On the contrary, it became his life's task
to work the two worlds together and to find the law of their connection. In
opposition to Romanticism's philosophy he would go from below upwards, not from
above downwards. And where empirical inquiry gives out, he would rather let
imagination range in the boldest analogies than let abstract thinking spin its
slender threads. His bent for paradox he gave vent to in a series of humorous
writings, under the pseudonym Dr.~Mises. His remaining writings divide naturally
into two groups. The one --- within which \textit{Zendavesta} (1851),
\textit{Ueber die Seelenfrage} (1861) and \textit{Die drei Motive des Glaubens}
(1863) are especially to be named --- develops his poetic-speculative world-view;
to the other belong his most epoch-making works in the fields of natural
philosophy and psychology: \textit{Ueber die physikalische und die philosophische
Atomenlehre} (1855), \textit{Elemente der Psychophysik} (1860) and
\textit{Vorschule der Aesthetik} (1876).

Right into his advanced old age Fechner was intellectually active and productive.
Even in his eighty-sixth year he made an interesting contribution to the
discussion of his psychophysical theory (\textit{Ueber die psychischen
Massprincipien und das Weber'sche Gesetz}, Wundt's Philosophische Studien IV).
% [p. 528]
Shortly afterwards he died, in 1887.\footnote{Interesting biographical notes in
A. Elsas: \textit{Zum Andenken Gustav Theodor Fechners} (Grensboten 1888). A fine
characterization of his great predecessor was given by Wundt in his memorial
words at the graveside (printed in ``Philosophische Studien'' IV) and later, more
fully, in his address on the centenary of Fechner's birth (\textit{Gustav Theodor
Fechner.} Rede. Leipzig 1901). A rich biography has been given by J. E. Kuntze
(\textit{G. Th. Fechner.} Ein deutsches Gelehrtenleben. Leipzig 1892), and Kurt
Lasswitz has given a critically searching presentation of Fechner's scientific
standpoint (\textit{G. Th. Fechner.} Stuttgart 1896, the first volume of
``Frommanns Klassiker der Philosophie'').}

\subsubsection*{α.\quad Poetic-Speculative World-View}

In the field of the world-view Fechner combats every conception which takes the
rich and bright world of mental life to be borne by, or to be the product of,
dark, intuitable things or beings. Materialism's ``matter'', spiritualism's
``soul-substance'' and Kant's ``Ding an sich'' are again and again objects of his
polemic. No less does he combat those who separate God and world, spirit and
nature. He rejects in equal measure belief in a natureless God and belief in a
spiritless matter --- orthodoxy as much as materialism. He reproaches the usual
conception of the world with dualistically separating the infinite and the
finite: ``One sets the infinite over against, above, beyond, outside, behind the
finite, perhaps even places an immense partition between them, as though they had
nothing whatever to do with one another. And yet they by no means lie outside one
another. On the contrary, the finite is the content of the infinite~.... The
infinite is therefore not incomprehensible either; it can be grasped by countless
handles within the finite; only it cannot be encompassed.'' (Ueber die
Seelenfrage, p.~111.)

Closely connected with the prejudices mentioned is, according to Fechner, the
opinion that only human beings and animals are ensouled, and nothing else --- at
least upon our earth. It is thought that experience forbids extending ensoulment
further. But only of my own soul have I direct experience; to the existence of
other souls I must infer by way of analogy. And what then can prevent me from
extending the analogy further, from human beings and animals to plants and to the
heavenly bodies, when good grounds show themselves for it? That plants have no
nerves is no counter-proof; for the lowest animal species have no nerves either.
If it is said that the plant is no closed individual, the answer is that
individuality is always relative, since no living being stands absolutely
separated from its surroundings. There is so continuous a transition between
animal and plant life that it is unwarranted to lay down an opposition between
them as between the ensouled and the soulless. The plant's mental life might
stand as far below the animal's as the animal's below man's. --- And why should
the heavenly bodies not be ensouled? Human beings and animals belong to the
earth, and the earth-soul might stand to the individual human and animal souls as
the earth's body does to the individual human and animal bodies. It is only an
artificial abstraction when one sets human and animal life as ensouled in
opposition to the life of the whole globe. The living does not proceed from the
lifeless; rather the lifeless is a product of organic metabolism and of the
perishing of organisms --- a conception which according to Fechner is easier to
defend than the opposite. In his work \textit{Ideen zur Schöpfungs- und
Entwickelungsgeschichte der Organismen} (1873) he develops especially the thought
that variation and differentiation proceeded with greater rapidity in earlier
periods of the earth than in the present one. --- The lower souls must be
supposed to stand to the higher as representations and motives stand to the
individual soul. In the end all souls are parts of the highest, all-embracing
soul, whose life and whose reality make themselves known in the law of causality
--- the principle of all particular laws of nature, and thereby of all connection
and order in existence.

As for Lotze, then, the law-governed connection in the world is the foundation
upon which the religious-philosophical view is built. What for many makes the
idea of God impossible or superfluous makes it, for Fechner and Lotze, precisely
necessary. In that universal law they find the highest unity expressed --- the
eternal and unchangeable which embraces all. The concept of law is the last thing
in all our knowledge; upon it, therefore, according to Fechner, our highest idea
must be built. (See especially his paper \textit{Ueber das Causalgesetz} in
``Berichte der sächs.\ Societät der Wissenschaften'' 1849. Cf.\ Seelenfrage
p.~205\,f.; Atomenlehre, 2.~Aufl., p.~125.)

The concept of the world goes back for Fechner, as for Lotze, into the concept of
God. The widest conception of the concept of God Fechner regards as the best
grounded; the narrower is in his eyes an abstraction. And since the world's life
is God's life, this cannot be complete in itself, but develops and unfolds with
and through the world's development. God's perfection consists not in finished
completeness but in unlimited advance. The divine world-plan is not finished once
for all, but is shaped in accordance with the world's events, which are at the
same time divine experiences. --- In this doctrine --- in which Fechner follows
his friend Weisse, while Lotze did not think he could ascribe validity to time
for the Godhead --- Fechner sees a carrying through of Christianity's declaration
that God is spirit and shall be worshipped in spirit and in truth, and that in
him we are, live and move: declarations which for the most part stand as mere
words! He admits that his doctrine is certainly not Christian, if one reckons as
an essential constituent of Christianity ``belief in the bite of the apple in
paradise and its mystical consequences, in the eternal damnation of the
non-elect, in miracles against the laws of nature, in God's remoteness from the
world, and in all the unedifying stuff out of which the theologians build up
Christianity'' (Seelenfrage p.~194). Fechner's bold theology makes it easier for
him to understand how man can be God's fellow-worker, and how evil can be foreign
to God's essence and yet a necessary member in the world's development.

\subsubsection*{β.\quad Psychophysics}

Early on, in accordance with his whole fundamental view, Fechner had arrived at
the conviction that the difference between the spiritual and the material could
not be a difference between two beings, of which the one was immaterial and the
other spiritless. The material world is the Godhead's outer, the spiritual world
the Godhead's inner side. The difference is phenomenal, resting upon a difference
in the spectator's standpoint, not upon a difference in substance. (Zendavesta
II, p.~341.) Or, as he later expressed it (Elemente der Psychophysik I, p.~2),
the difference between the spiritual and the material is to be likened to the
difference between the concave and the convex side of one and the same circle: he
who stands within the circle sees only the concave side, he who stands outside it
only the convex --- and he who can change his standpoint may perhaps believe that
he has two different things before him. It is Spinoza's identity hypothesis that
Fechner here arrives at, led by the physicist's interest in maintaining the
continuity of the physical processes, and by the philosopher of religion's
interest in thinking the Godhead, the all-embracing spiritual principle, in such
a way that it expresses itself at every point in material phenomena, so that
these are not loosened from it and do not fall outside it. Fechner arrives, so to
speak, at a theophysics before he develops a psychophysics.

A confirmation of this fundamental conception he found in the law of the
conservation of energy, and he is the first to use the newly discovered law in
this way. (Elemente der Psychophysik I, p.~21--45.) He admits that no proof has
been given that this law holds for the material processes bound up with spiritual
activities; but he holds that all experience points in this direction, and that
we must therefore, until a counter-proof is produced, assume the law to be valid.
Experience establishes that the material processes to which conscious life is
bound stand in interaction with the other material processes within and outside
the body, so that the total physical energy at our disposal may now be consumed
predominantly by the former, now by the latter. Physical energy is drawn upon
both when we think and when we chop wood --- and therefore we cannot do both at
once as well as we can do each of them separately.

% [p. 532]
From Spinoza he differs in this doctrine only by laying strong weight upon
finding out a mathematical functional relation\footnote{In my earlier
presentation (\textit{Filosofien i Tyskland.} 1872, p.~264\,f.) I allowed myself
to be misled by some utterances in \textit{Zendavesta} (II, p.~352\,ff.) into
saying that Fechner assumes an interaction between spirit and matter. What
Fechner there develops is, however, only the importance of changing standpoint
and of finding out the formula for the mutual relation of the two sides. As
little there as later does he assume a passage from the one side to the other. In
\textit{Elemente der Psychophysik} I, p.~8 it runs: ``\textit{Wir nennen das
Psychische Function des Physischen, davon abhängig und umgekehrt, insofern eine
derartige constante oder gesetzliche Beziehung zwischen beiden besteht, dass von
dem Dasein und den Veränderungen des Einen auf die des Anderen geschlossen werden
kann.}'' For this kind of dependence Fechner often uses (for example Elemente I,
p.~18; II, p.~393) the expressions \textit{Simultanbedingungen, simultane oder
Wechselabhängigkeit}. On his conception the brain processes corresponding to the
activities of consciousness cannot --- as Lotze and the spiritualists hold --- be
regarded as releasing impressions (\textit{Reize}) upon the soul (Elemente I,
p.~18).} between the two sides of existence. He first proceeded from the view
that there must be a directly proportional relation between them. Later --- he
has noted the date and the situation: 22 October 1850, in the morning, in bed ---
he hit upon the thought that the spiritual did not rise and fall directly with
the material, but that the changes of the former corresponded to the
\emph{proportional} changes of the latter, so that the change ($d\gamma$) in the
intensity of a spiritual state is to be set equal to the ratio between the change
($d\beta$) in the energy of the corresponding material state and the energy
($\beta$) already present --- that is, $d\gamma = d\beta/\beta$. With this he had
come down from the region of indeterminate speculations to an assumption that
could be checked experimentally. He found it at once confirmed by the fact that
the strength of the sensation of light does not grow as fast as the physical
intensity of the light. His own experiments and his researches in the literature
gave him a material on the basis of which he set up the general law that a
sensation's increase does not correspond directly to the increase of the physical
impression, but to the ratio between that increase and the whole impression
present --- so that with one and the same increase of impression the sensation
grows the less, the stronger the impression already is. This law Fechner calls
the Weber law, in memory of his teacher, the physiologist E. H. Weber, from whom
some of the most significant experiments upon which it was built derived. --- Of
the law's validity, its limitation and its interpretation nothing further shall
be said here; for that, reference must be made to more recent psychological
works. Here it need only be remarked that by the step Fechner thus took, exact
experiment became possible in the psychological field: he made --- to use an
expression of Galilei's --- something measurable that had not hitherto been so.
Through its relation to the outer impression the sensation, which in and for
itself, as a purely subjective element, cannot be measured, can be determined
quantitatively. Therewith the road was pointed out which \emph{experimental
psychology} had to take: it had to seek points of connection between psychical
processes and such processes as can be directly counted, weighed or measured.

What Fechner had first been seeking was certainly not what the new science came
to occupy itself with. He had wanted to find the relation between the mental
activities and the corresponding brain processes. But that relation is
inaccessible to us. What we immediately apprehend is, from the subjective side,
the mental state, and from the objective side, the physical impression. For the
relation between these the law Fechner found holds, within certain limits; and it
is most probable that brain process and mental process are directly proportional
--- which also agrees best with the identity hypothesis Fechner espoused.
Psychophysics has so far become something other than Fechner himself had
conceived. But this does not lessen the importance of his brilliant work
\textit{Elemente der Psychophysik}: between the human sciences and natural
science there has been inserted --- in consequence of the great law of the
division of labour --- a new discipline. And it was in Fechner's free spirit,
rich in possibilities and in ideas, that this budding was accomplished. The
investigator who after Fechner has done most in the new science's field said at
his graveside: ``Psychophysics, which he founded, was only the first conquest in
a domain whose further appropriation could no longer present significant
difficulties, once this beginning had been made'' (\textsc{Wilhelm Wundt}:
\textit{Zur Erinnerung an Fechner.} Philos.\ Studien IV, p.~477).

In defence of the psychophysical law he had set up, Fechner composed a series of
writings which are models of truth-loving and careful discussion and of
chivalrous polemic; they are marked too by the author's freshness and humour. ---

Akin to Fechner's endeavours towards the founding of an experimental psychology
is his endeavour towards an aesthetics grounded upon experience. His
\textit{Vorschule der Aesthetik} (1876) is an epoch-making work, into whose
special content we cannot here enter more closely.

\subsubsection*{γ.\quad Philosophy of Nature}

His opposition to Romanticism's philosophy Fechner has expressed above all in his
``Atomenlehre''. Speculation, he says, descends into nature as the bear into a
beehive. It enjoys the coherent honey, but does not notice that it has been
brought together by the activity of a multitude of small individual beings.
Natural science shows --- in physics and in chemistry alike --- that the material
processes can be understood only by being regarded as processes of interaction of
small parts which for us are indivisible. Every part of the material world,
however small, natural science conceives as a little world with globes and
systems of globes. What seems to us a continuous mass is, like so many nebulous
masses in the sky, a group of different parts. Natural science does not assume
absolute atoms, but only that divisibility goes further than the eye and the
microscope discover. It assumes only relative atoms --- that is, atoms
indissoluble for us, or by the forces of nature known to us --- led by the need
to refer every effect to a determinate place in space as its point of departure.
Philosophical atomism goes to the end of the road along which physical atomism
had reason to travel only part of the way. We obtain a philosophical conclusion
to the conception of nature when we assume that the smaller one thinks the atoms,
the more exact will be the results one can derive by thinking of matter as
consisting of atoms. It was in order to have points of departure for the effects
that natural science conceived atoms. The share with which every such point of
departure contributes to the fulfilment of a law of nature we call the force
proceeding from that point. What we know about the atoms is now precisely only
that these contributions proceed from them: therefore philosophical atomism
designates them as centres of force. There is thus no reason to ascribe extension
to the atoms --- the absolute, philosophical atoms --- and the concept of matter
is no longer materialistic. Philosophy stands here for Fechner as a hypothesis in
which we think as absolute what in reality is always given only relatively.

The foundation of our whole knowledge of existence is the law of the connection
of phenomena; from the law we come to determine the forces in the co-operating
elements. Fechner will not, with Leibniz, Herbart and Lotze, see psychical
elements in the atoms. These are for him merely the last points at which
spiritual beings arrive when they analyse the content of their consciousness.

A closer explanation Fechner regards as superfluous. The atom is the lower limit
of our knowledge, just as the universal world-law, which testifies to the reality
of an all-embracing being, is the higher limit. Between these two limiting
concepts lies all our knowledge. (Seelenfrage, p.~215--216.)

\subsection*{c.\quad Eduard von Hartmann}

In the year 1869 there appeared a book that caused a great stir and enjoyed a
circulation rare for voluminous philosophical works. It was Eduard von Hartmann's
\textit{Philosophie des Unbewussten}. Its motto --- ``Speculative Resultate nach
inductiv-wissenschaftlicher Methode'' --- shows that the author aimed at a
grounding and development of philosophical ideas similar to Lotze's and
Fechner's. By his clear and broad presentation, by his detailed use of
natural-scientific examples, by meeting the reader at once with a realist
foreground and a romantic-mystical background, and finally by taking up pessimism
--- which, under the influence partly of Schopenhauer's now much-read writings,
partly of prevailing currents of the time, was a widespread train of thought ---
as a member of his system, Hartmann could find readers in wide circles. The
attention he aroused is seen from the fact that in the course of twenty years ten
editions of his principal work appeared, and that in the years 1870--1875 alone
there appeared fifty-eight writings dealing with him.

\textsc{Eduard von Hartmann} was born in Berlin in 1842, was the son of a
general, and himself entered upon a military career. In his free time he
cultivated music, painting and philosophy. A knee injury compelled him in 1865 to
resign his commission. And when he became clear that his real talent did not lie
in the field of art, he gave himself --- ``bankrupt in everything but thought''
--- to philosophy. He had philosophized and worked out philosophical essays
earlier. He now proceeded to the working out of ``Philosophie des Unbewussten'',
which must still be regarded as his principal work, notwithstanding that it has
been followed by some thirty larger and smaller works. The most considerable of
these is certainly \textit{Phänomenologie des sittlichen Bewusstseins} (1879),
which is on the whole Hartmann's best work. For the characterization of his
philosophical standpoint I shall here keep to two principal points, in full
consciousness that Hartmann's untiring activity in developing his ideas in all
directions, and illuminating them from all points of view, does not thereby quite
come into its own. Hartmann has complained that I have ``halved'' him by
mentioning only the part of his activity that lies before 1880. That, however,
lay in this book's plan; and I have besides, it must be admitted, the view that
the two works named are and remain Hartmann's most considerable contributions to
philosophy. Those who wish to supplement the short presentation of Hartmann's
philosophy that I can give here I must refer to the detailed work, sprung from
excessive admiration, of \textsc{Arthur Drews}: \textit{Hartmanns philosophisches
System im Grundriss} (Heidelberg 1902).

\subsubsection*{α.\quad Philosophy of Nature and Psychology}

While Lotze and Fechner in principle laid the natural-scientific mode of
explanation at the ground, and would ground their philosophy by drawing the
consequences of its presuppositions and results, Hartmann's philosophy bears far
more the stamp of a neo-Romantic reaction against natural science's realism. He
sets out to show that the natural-scientific mode of explanation is not
sufficient, but that together with the causes the mechanical view of nature
assumes there must operate a spiritual principle which he --- in order not to
make it too anthropomorphic --- calls \emph{the Unconscious}. He appeals to this
principle wherever he thinks the causes empirical science can exhibit do not
suffice. It is thus that he thinks to arrive by ``inductive'' means at
``speculative'' results.

Hartmann agrees with Fechner that matter must be conceived as a system of atomic
forces. No single atom can have mass, any more than one can speak of the
magnitude of a unit. If, then, we keep to atomic force as the last element to
which the analysis of material reality leads, there can --- Hartmann holds --- be
determinate connection and consistency in the atomic force's striving only when
that striving is conceived as a willing, with an unconscious representation of
its goal. We understand force only by thinking of it as a willing. Matter itself
is therefore representation and will --- and the difference between matter and
spirit falls away. In organic growth there likewise expresses itself an
unconscious willing and representing: in the organism's coming to be a purpose is
realized, a purpose in relation to which the single substances and processes are
to be regarded as means, whose provision and combination become intelligible only
through the purpose, notwithstanding that this is not an object of consciousness.
There is only a difference of degree between organic formation and instinctive
action. Instinct cannot, Hartmann claims, be explained by the material
organization or rest upon a nerve-mechanism implanted once for all; for then it
could not vary in its expressions when the organization does not change. As
little does it rest upon conscious deliberation. There remains, then, only to
conceive it as unconscious willing and representing. In the human mind the
Unconscious expresses itself in sense-perception, in that the intuition of the
outer object comes to be by an unconscious working-together of the sensations.
The association of representations proceeds only through representations being
brought forward, without conscious search, which stand in connection with the
present ones. Feelings and motives arise on account of unconscious processes and
are therefore often unintelligible to ourselves. Even where the will is bound up
with consciousness, an unconscious will is required if it is to lead to a willed
movement: consciousness does not know which nerve-centre in the brain must be
affected for the movement to come about. In history the Unconscious is able to
get individuals to work in the service of great world-purposes while they believe
they are working only for their own limited ends. What is fate or providence
other than the Unconscious's dominion in men's actions, so long as their conscious
understanding is not mature enough to make history's purposes its own? --- The
Unconscious is everywhere that which makes co-operation between individual beings
possible. There are individuals of every possible degree --- from the atom to
universal nature --- and it is one and the same Unconscious that stirs in them
all. Plurality has only phenomenal, not metaphysical significance. Hartmann
declares himself in principle in agreement with speculative theism (as it appears
in Schelling, Weisse, Lotze and Fechner), his ``Unconscious'' being rather to be
called ``superconscious''. But he prefers the negative expression in order to
keep anthropomorphic representations at a distance. He thereby overlooks that
even if the name he has chosen removes such representations, they are in high
degree favoured by the manner in which he lets ``the Unconscious'' operate.

It is an unscientific procedure that Hartmann employs when, wherever he thinks he
finds a hole in the scientific explanation, he puts his ``Unconscious'' in as a
magical means that works equally well everywhere. Of it holds what Galilei said
of the appeal to an almighty will: it explains nothing, because it explains
everything equally well. The real interest in the concept of the Unconscious has
in Hartmann scarcely arisen from a striving after understanding of the phenomena
either. It is due rather to a need for rest from reflection, analysis and
criticism, and to a faith in the value of what emerges involuntarily and
unprepared. Thus it runs in one place in Hartmann (Philosophie des Unbewussten,
3.~Aufl., p.~369): ``Conscious reason is only negating, criticizing, controlling,
correcting, measuring, comparing, combining, subordinating, inducing and
generalizing. Never is it creatively productive, never inventive; in this respect
man is entirely dependent upon the Unconscious, and when he loses the capacity to
receive the Unconscious's promptings he loses his source of life, without which
he will drag out his life in monotonous fashion under dry schematizing. Therefore
the Unconscious is indispensable, and woe to the age that violently suppresses it
and, in one-sided overvaluation of the consciously rational, will let that alone
count!'' Consciousness cannot be the purpose of development, but can only be a
means, necessary at a certain stage of development. (Phil.\ d.\ Unbew.,
3.~Aufl., p.~739.) --- It is this psychological-historical experience of the
significance of unconscious life that is the warranted foundation of Hartmann's
thinking. But he then makes the Unconscious into a mythological being, a demon
that everywhere intervenes in the natural connection, directs the atoms,
displaces the brain molecules, adjusts the circumstances. Hartmann's philosophy
takes on a mythological and dualistic character, which is certainly most
prominent in the first editions of ``Philosophie des Unbewussten'', but is too
closely bound up with his principle for it to be able, despite all correctives,
to disappear entirely. Characteristic above all is his position on Darwinism, in
that he seeks to show that the selection of qualities does not explain the
arising of new forms, but is only a mechanical means which ``the Unconscious''
makes use of when it wants to prepare the way for such forms. Darwin keeps ---
Hartmann holds --- to the conditions, but overlooks the properly producing force.

It was then no wonder that Hartmann was often attacked from the natural-scientific
side. Now, in order to show that he himself had a very good command of the
natural-scientific points of view, he wrote an anonymous critique of himself:
\textit{Das Unbewusste vom Standpunkt der Physiologie und der Descendenzlehre}
(1872). The critique was so good and so telling that his opponents praised the
anonymous work and used it against him. In the second edition Hartmann named
himself and added notes in which he sought to refute the objections he had
himself raised. It is doubtful, however, how far he has succeeded; he has
conjured up spirits that will not vanish again. He admits that he had undervalued
the mechanical causes and overlooked many intermediate links. But his procedure
continues to be that of seeking out the holes which scientific inquiry always
leaves, and filling them by the help of his mystical \textit{deus ex machina}.
Since there is always the possibility that the holes may be filled
scientifically, it is a dubious method to employ. Lotze employed this method in
his spiritualistic psychology; Hartmann extends it to his whole philosophy,
without attending to Lotze's clear and energetic emphasis upon the significance
of mechanism. His ``speculative results'' are not won by ``inductive method''. He
accordingly admits (Das Unbewusste etc., 2.~Aufl., p.~359\,f.) the insufficiency
of induction, and ascribes to it only the significance of serving to confirm what
has been thought out by another way.

\subsubsection*{β.\quad Pessimism and Ethics}

Hartmann declares it a great misunderstanding that he should himself, in his
personal temper, be a pessimist; he teaches pessimism as a \emph{theoretical
view}, after disappointments and sufferings in his youth opened his eyes to it.
``With the chapter on pessimism [in ``Philosophie des Unbewussten''] I have'', he
says, ``written myself free once for all from world-sorrow as such, and purified
it into an objective, calm \emph{knowledge} of existence's wretchedness --- but
thereby also won back for myself that undisturbed cheerfulness which may be found
in the philosopher who floats in the ether of pure thought and from there regards
the world, and his own pain, as an alien object of investigation.''
(\textit{Mein Entwickelungsgang.} Gesammelte Studien und Aufsätze, p.~38.) He
emphasizes as well that pessimism does not even theoretically constitute his
whole philosophy, and does not give it its principal stamp. (\textit{Die Stellung
des Pessimismus in meinem philosophischen System.} Zur Geschichte und Begründung
des Pessimismus, 2.~Aufl., p.~18--28.) He describes his task as that of uniting
pessimism --- as the result of valuing the world by the standard of happiness ---
with the philosophy of development. He would thus reconcile Schopenhauer's
philosophy with Hegel's, just as in his very concept of the Unconscious he has
sought to combine Schopenhauer's absolute will with Hegel's absolute Idea. In a
peculiar way Hartmann here finds a point of attachment in Kant. When Kant
declared every attempt to solve the problem of evil hopeless, that was consistent
from his standpoint, since he proceeded from the Christian concept of God, with
which the reality of evil stands in insuperable contradiction. But the problem
must be posed thus: how must the world's absolute ground be thought, if it is,
without contradicting itself, to have come to posit such a world? And when the
question is so posed, Kant has prepared the way for its treatment, since he has
given indications in a pessimistic and in an evolutionist direction alike.
Schopenhauer pursued the one indication, Hegel the other: now the time has come
to gather up what lies in both. (\textit{Kant als Vater des modernen
Pessimismus.} Zur Gesch.\ u.\ Begr.\ des Pess., 2.~Aufl., p.~136\,f.) (Compare
Note~17 in the present volume.)

Since now --- so Hartmann reasons --- an unprejudiced consideration shows more
unhappiness than happiness in the world, the world's coming to be cannot be
derived from reason. It cannot be due to a rational but must be due to an
irrational principle. The explanation must then be sought in the will-element
within the Unconscious having separated itself from the representation-element.
Hartmann takes up Böhme's and Schelling's mystical doctrine of a dark element in
the Godhead whose breaking loose leads to the world's development. In a
groundless manner the in itself blind willing has torn itself out of connection
with representation, or the Idea. But this seeing element in the Unconscious
continues to work in the world's development and seeks to lead that blind and
refractory element to harmony and reconciliation. Hence the double-sided
impression we get of the world. Both pessimism and evolutionism are right. There
are two principles at work in the world. Hartmann's philosophy of religion comes
here to recall Voltaire's, Rousseau's and Stuart Mill's --- save that development
is, on his belief, a progressive redemption of the suffering God, since what
matters is, through the abolition of the will, to undo what was done when the
world became actual. Of great importance in this respect is that the knowledge of
life's wretchedness should awaken, and Hartmann therefore holds that with
Schopenhauer a new world-period has begun. (Phänomenologie des sittlichen
Bewusstseins, p.~782.) Only slowly does this knowledge dawn in mankind. It
happens through a successive abolition of illusions. First one expects happiness
in the present life. When one discovers that the goods it can bring are only
apparent, one sets one's hope upon another life. That is the illusion's second
stage. Belief in the beyond cannot maintain itself; one gains fresh confidence in
earthly life, but, taught by experience, one does not expect an immediate
enjoyment of happiness, and instead hopes for a state of bliss for mankind at
some time in the future, straining every nerve to prepare the way for it. That is
the illusion's third and last stage, at which mankind will probably remain for
long ages. But in the very fact that a small minority has seen through the
illusion in all three of its forms, Hartmann sees sufficient warrant for the
pessimistic theory; its correctness cannot be settled by a vote.

Hartmann's ethics hangs closely together with his pessimistic theory. There is
for him an irreconcilable conflict between culture and happiness. Culture's
advances are bound up with a decline in happiness. The more complicated life's
apparatus becomes, the more occasion for displeasure is given. Susceptibility to
pain grows greater, and the increasing reflection sees through the illusions more
easily. Culture lets needs grow faster than the means of satisfying them. One
must then choose between culture and happiness, between the doctrine of
development and the doctrine of bliss. Happiness presupposes peace and rest ---
and will therefore bring stagnation and dissolution with it. Development leads
ever further, until all possibilities are
% [p. 542]
exhausted.\footnote{On the relation between culture and happiness compare my
\textit{Etik}, Kap.~7, and my \textit{Etiske Undersøgelser}, Kap.~2.} Through
development dissatisfaction is increased and the illusions are abolished, until
life at last comes to a stop. Then redemption is reached, and the blind will's
breaking loose is annulled. The highest form of ethics is according to Hartmann
that which builds not merely upon love of mankind and the urge to development,
but makes compassion with God --- the need to release the bound and suffering
world-principle --- the foundation of morality. One then makes the Unconscious's
purposes one's own.

Hartmann's philosophy is theology inverted. He operates confidently with the
thought of an absolute world-purpose, without seriously engaging the difficulties
that thought contains. In his ethical work (Phän.\ d.\ sittl.\ Bew.) there are,
however, excellent particulars, especially in the characterization and criticism
of the different standpoints. And in detail his train of thought is independent
of pessimism's metaphysics. He has overcome that metaphysics to a still higher
degree than he says. In practice he is an evolutionist and an optimist. For
provisionally the affirmation of the will to life is the one right thing, since
something can be accomplished for the world-process only by full devotion to life
and its pains, not by cowardly renunciation and self-surrender. Pessimism's time
really comes only after long ages --- ages so long that one wonders doubly at
Hartmann's confident knowledge of what will happen when they are past. By this
long postponement pessimism loses the personal, practical significance it had in
Schopenhauer. Hartmann stands to Schopenhauer as the modern Christians, for whom
the day of judgment has been pushed out into the indefinite, stand to the first
Christians, whose whole life was determined by the expectation of its speedy
coming. Distant possibilities cannot be determining for our conduct of life, and
cannot shake life's value. Only the life we know and live can we value, and
within that life itself we must find what we are to be able to acknowledge as
valuable.

\section*{4.\quad Criticism and Positivism}
\label{sec:kriticisme}

\subsection*{a.\quad Friedrich Albert Lange}

It was natural that the dissolution of Romanticism's philosophy, the rise of
natural science, and materialism's claim to be science's last word, should bring
to greater influence the tendency which in the century's first half had been only
an undercurrent. Already the greater attention and adherence which Schopenhauer's
and Herbart's ideas now became the object of are signs in this direction. There is
a need for a critical revision of the very foundation of our knowledge, in order
to be able to take up the right attitude to the new material of knowledge. In
Lotze and Fechner, and in part in Hartmann too, investigations in this direction
are indeed to be found, but the problem of knowledge proper is not taken up. They
take realism as given, and then seek, more or less consistently, to find a
foundation for their idealist constructions either in realism's presuppositions
or in its holes. Great as the interest of Lotze's attempt in this direction is,
it was nevertheless the problem of existence rather than the problem of knowledge
that he treated. One of the first to draw attention to the necessity of taking up
Kant's work upon the problem of knowledge again was the famous historian of
philosophy \textsc{Eduard Zeller}, in a small work \textit{Ueber die Bedeutung
und Aufgabe der Erkenntnisstheorie} (1862), which caused all the more stir in
that its author had come from Hegel's school. A few years later
\textsc{Friedrich Albert Lange}, in his brilliant historical-philosophical work
\textit{Geschichte des Materialismus} (1866), pointed to epistemology as a court
of appeal overlooked not only by Romantic but also by materialist philosophy, and
sought upon its basis to carry on the struggle for an idealist world-view.

Lange is one of the most attractive figures in the history of recent German
philosophy. Ideal enthusiasm, sharp criticism, great fullness of knowledge and a
vigorous realist eye have entered in him upon a rare union. He commands besides a
great talent as a writer. Unfortunately his career was short. He was born in 1828
near Solingen in Rhenish Prussia, but lived his early youth in Zürich, whither
his father, the well-known theologian J. P. Lange, was called as professor of
theology in 1841. He thus received the foundation of his intellectual development
in the free air of Switzerland, and it was to be traced through his whole life
and striving. His university studies he pursued in Bonn. Educational work stood
before him as his principal aim, and he studied philology as a means to it,
without however neglecting natural-scientific and philosophical studies. His
philosophical direction was early laid down. In a letter of 1851 (in Ellissen:
\textit{Friedrich Albert Lange.} Eine Lebensbeschreibung. Leipzig 1891,
p.~69--73) he develops the view that all representations man can form of
existence --- of God and the world, of good and evil and so on --- are developed
according to the laws of the human mind: those laws are then the last foundation
of our knowledge, beyond which we cannot get. He was here led to the
psychological-epistemological standpoint which Ludwig Feuerbach had reached
through criticism of the religious and speculative ideas, though Feuerbach does
not seem to have influenced him. What influenced him philosophically in his youth
was above all Hegel's ``Phenomenology'', Herbart's psychology and
Schleiermacher's writings --- of the last he speaks with enthusiasm in a letter
of 1849.

After some time as lecturer and schoolmaster, Lange worked as editor and as a
political and social agitator. He took part in the opposition to Bismarck's
domestic policy and in the German workers' associations' efforts to win a firm
organization. The work \textit{Die Arbeiterfrage} (1865) derives from this
period. It is one of the first writings --- next after Mill's political economy
--- in which the social question is conceived from a more general point of view
and in an impartial manner. It has interest besides in bringing the Darwinian
ideas into connection with the workers' question, in that this is seen to spring
from the struggle for existence in the social field. Lange conceives the matter
thus, that what is required is to wage a struggle against the struggle for
existence --- a struggle which precisely constitutes man's higher spiritual
vocation. (Since at this time only ``The Origin of Species'', and not ``The
Descent of Man'', had appeared, Lange could not make use of the points of view
Darwin has given for conceiving the ethical endeavours too as a struggle for
existence.) He saw the political significance of the workers' question in its
exerting a constant pressure upon the conservative institutions, whose dykes will
be broken through if canals are not dug in time. The liberal-minded men in the
higher classes now have the choice between strengthening those dykes and helping
to dig the canals. Lange's own choice was easily made. But in Germany men were
busiest with the dykes, and Lange therefore emigrated to Switzerland. In the
midst of his journalistic and agitational activity he had nevertheless found time
to put the last hand to his principal work \textit{Geschichte des Materialismus}
(which according to Ellissen, p.~145, appeared in the autumn of 1865, though the
title page bears 1866). It is not a purely historical work. Through a spirited
depiction of materialism's principal forms and of the development of the natural
sciences, Lange seeks to illuminate the problem of the relation of the spiritual
and the ideal to the material, and he sets the theoretical question at the same
time in connection with the practical and social one. Although many parts of the
work are now of course out of date, no one will read it without enrichment and
stimulus. --- Shortly afterwards there appeared a new social-philosophical work
of Lange's: \textit{Mill's Ansichten über die soziale Frage} (1866).

In Switzerland Lange worked as bookseller, teacher and editor in Winterthur, and
unfolded an astonishingly many-sided activity. He took eager part in the
agitation for a revision of the constitution of the canton of Zürich. Later he
became professor in Zürich and could now live for his studies again. He was soon
noticed, however, in his old fatherland, back to which he also longed, and in
1872 he became professor in Marburg. But by then he had already suffered for
several years from cancer of the abdomen, and his last years were painful. His
ideal disposition and his great enthusiasm for research and for freedom sustained
him; his spiritual nourishment he found above all in Schiller's philosophical
poems. He died in 1875. --- An unfinished but interesting work on logic
(\textit{Logische Studien.} 1877) was published shortly after his death. Later
(1897) Ellissen published \textit{Einleitung und Kommentar zu Schillers
Philosophischen Gedichten}, a work Lange had prepared during his stay in Zürich
but laid aside for other labours. Schiller was, as Lange said in the preface to
his principal work, the Kantian with whom he felt himself most akin. Schiller's
poems contained Lange's religion.

As the title of Lange's book indicates, his problem is how one is to take up a
position with respect to materialism. His standpoint here is the same as
Fechner's. He would overcome materialism precisely by carrying it through. What
he praises in materialism, and what for him constitutes its great historical
significance, is its consistent demand for the exhibition of material causes for
all natural phenomena. In idealism, by contrast, Lange sees a great danger in the
theoretical field, since it easily leads to relaxing the strict mechanical
explanation and to inserting subjectively constructed members into the series of
objective causes. The functions of the nervous system and of the brain too must
be explained according to the general physical and chemical laws. One must ---
relying upon the law of the conservation of matter and force --- follow the
influences from without through sense-organs, afferent nerves, brain, motor
nerves and muscles, and seek a complete and unbroken connection in this whole
process. It is unwarranted to assume the connection broken at any point.
Consciousness itself is not a member of this connection. It is the subjective
state of the individual in whom the process takes place --- another side of the
process, but not itself a part of it. Here materialism has its limit: precisely
because the states of consciousness are not members or parts of the material
processes, and cannot be explained by the law of the conservation of matter and
force, precisely therefore do they fall outside materialism, and materialism
comes into contradiction with itself when it thinks it has explained everything.
It can at most explain the objective, material process, which is the way in which
the conscious individual's subjective state appears to outward consideration.
Lange regards it as conceivable that one will some day be able to determine that
part of the material processes which coincides in time with a certain state of
consciousness in the individual; but he does not think one will ever be able to
give a closer determination of the relation between the subjective state of
consciousness and the objective nerve process. It is a warranted conjecture that
behind the two corresponding worlds --- the world of matter and the world of
consciousness --- there lies an unknown third as their common cause.

It is remarkable to see how the great hypotheses of the seventeenth century are
renewed after the rise of natural science at the middle of the nineteenth. As
Vogt, Moleschott and Büchner renew Hobbes's standpoint, so Lotze renews
Descartes's, and Fechner and Lange Spinoza's. It is a testimony that these
hypotheses indicate the possibilities between which our thinking must here make
its choice. On closer investigation one will find that it is not a simple
resumption but a renewal of the older hypotheses that has taken place. Not only
the natural-scientific foundation but also the philosophical elaboration is other
than it was in the great dogmatists of the seventeenth century. There is greater
critical circumspection, a clearer consciousness that what is at issue is finding
the best and most consistent points of view in a discussion that can scarcely be
brought to a close --- not the setting up of definitive systems. One notices that
the critical philosophy, whose breakthrough lies between the two epochs, has not
worked in vain.

Lange's standpoint, which from one side points back to Spinoza, now precisely
from another side points back to Kant. He does not stop at the double form of
existence --- as motion and as consciousness --- but calls attention to the fact
that the whole of matter, the world outside us as well as our own organism, brain
and sense-organs included, is only there for our consciousness, is a product of
consciousness. It is in virtue of our organization that we apprehend the world as
we do apprehend it. That Kant showed, and the later physiology of the senses has
confirmed it. Here again materialism is overcome by being consistently carried
through. For even if it be conceded to materialism that our whole conception of
the world is a product of our material organization, that material organization
is itself only an object of consciousness, and exists immediately only in and for
consciousness. Matter and organization are our representations --- even if they
are representations we necessarily form, and form according to determinate laws.
The dispute between matter and spirit is thus ended in the latter's favour. The
naive belief in matter must vanish upon closer investigation, and it can vanish
without natural research's campaign of conquest being in the least degree
hindered thereby. --- In the way Lange here takes up his position against
materialism he has a predecessor in Schopenhauer. At about the same time as
Lange, the famous Viennese pathologist \emph{Rokitansky} expressed a similar
% [p. 548]
view.\footnote{Karl Rokitansky's spirited paper \textit{Der selbständige Werth
des Wissens} (Sitzungsberichte der Wienerakademie 1867) did indeed appear more
than a year after Gesch.\ d.\ Mat., but as early as 1862 he had in an address
pointed to Kant's idealism as the true consequence of natural-scientific
thinking. Compare Th. Meynert: \textit{Karl Rokitansky.} Ein Nachruf. (Sammlung
von populär-wissenschaftlichen Vorträgen. Wien und Leipzig 1892, p.~76\,f.) ---
Helmholtz too (in his \textit{Physiologische Optik.} 1867, p.~443\,ff.) maintains
that it is in the last instance unwarranted to ascribe to our representations any
validity other than a practical and symbolic one. The physiologist A. Fick
arrives at a similar result in his lecture \textit{Die Welt als Vorstellung.}
Würzburg 1870.}

% [p. 549]
There is in Lange's presentation in ``Geschichte des Materialismus'' a certain
wavering as to how far he will extend subjectivism.\footnote{Max Heinze has shown
this in his interesting paper \textit{Der Idealismus Friedrich Albert Langes}
(Vierteljahrschrift für wissensch.\ Philos.\ I, p.~183--185).} Occasionally he
says that the very opposition between the thing in itself and the phenomenon is a
product of our organization, and that we therefore do not know whether that
opposition has any significance whatever outside our experience. In other places
--- especially in the second edition --- he speaks of an order of things
corresponding to our knowledge. In a letter to the natural scientist Dohrn (in
Ellissen, p.~258--262) Lange develops the view that with respect to this whole
question four different stages must be distinguished: 1) the naive belief in
sensible appearance; 2) the recognition that the senses do not give us things as
they are in themselves, but vibrations of air as sound, vibrations of ether as
colour and so on --- that is natural science's standpoint; 3) the recognition
that our understanding as well as our senses belongs to our organization, and
that therefore even the scientifically corrected apprehension of sense does not
give us the thing in itself --- that is Kant's standpoint; 4) finally the
recognition that the very opposition between thing in itself and phenomenon is a
product of our organization, so that its validity may be doubted. This
fundamental doubt cannot, according to Lange, be removed by way of knowledge. It
marks science's limit, beyond which only poetry, not thinking, can lead. Lange
reproaches Kant for stopping at the third stage; Kant's Romantic successors he
reproaches with having ``tumbled out beyond the fourth stage into poetry, and
mixed criticism and poetry in their systems instead of bringing them into a
clearly delimited relation to one another''.

In this conception Lange has, however, overlooked that when our knowledge applies
at its very limit the same principles as within the limit, it must necessarily
hold fast to the opposition between subject and object, knowledge and thing: in
the contrary case it would in one form or another come to assume a creation out
of nothing, a causality proceeding from an absolute unity. It is a law for all
knowledge that what we understand is understood only by the help of \emph{several}
premisses or conditions.\footnote{Compare my \textit{Psykologi} V~D, and
\textit{Filosofiske Problemer} (Universitetsprogram 1902), p.~48--52.} If one
does not acknowledge that law, one becomes entangled in difficulties which from
Nicolaus Cusanus onwards emerge again and again in the history of thought.

Lange, like Schopenhauer, uses idealism --- in the sense of subjectivism --- to
unseat materialism. But he also espouses idealism in the word's practical sense.
It is not enough for him to maintain reality's subjective character; reality
would not thereby become other than it is, bad or good. But he maintains that
just as deeply grounded in human nature as the need to apprehend and explain
reality according to natural-scientific laws is the need to form ideal images and
to regard them as the expression of the highest reality. The producing power
which expresses itself in sensation and in thinking does not rest content with
what is given, but strives beyond it in order to form an ideal world. Man cannot
do without a supplementing of reality. All religions and all speculative systems
are products of this need. They differ from the scientific picture of reality in
that it is not the organization common to all, the race's nature, that operates
in their arising, but individual and personal tendencies and requirements --- and
also in that a complete confirmation by experience is not possible. It is of
great importance to hold fast to this difference, so that one does not lay down
as literal truth what is only a symbol. Every significance for the knowledge of
reality must be denied to the religious and speculative ideas. But on the other
hand a distinction must also be drawn between such ideas and mere figments.
Although Lange in single places lets poetry begin where knowledge ends, he
maintains in other places (1st ed.\ p.~539, 541; 2nd ed.\ II, p.~540) that
speculation stands between empirical science and poetry: in the one the material
rules, in the other free form; speculation seeks to unite the two. Clearer is the
difference Lange finds in this, that the idea is to be judged by its value for
the development of our spiritual life, not by its grounding and origin (1st ed.\
p.~346\,f.; 2nd ed.\ II, p.~545). We do after all regard Cologne cathedral as
something other and more than a heap of stones. Indeed, Lange even speaks of
practical idealism as possibly hanging together with the unknown truth, though in
another way than materialism does: the idea is image and symbol for an
unknowable Absolute. (1st ed.\ p.~346, 539, 541, 545. --- 2nd ed.\ II, p.~494
and elsewhere.)

Here again it shows itself how difficult it is to place an absolute limit. Lange
becomes here involuntarily a Platonist. How does he know that the ideas which for
us express the highest value are images or symbols? And on the other hand --- if
his sharp distinction between ideas and scientific knowledge is held fast --- who
will then work at the development of the ideas, if they have no real significance
at all? The speculative ideas can be maintained in their significance only by
being seen as hypotheses, in which human thought draws the last consequences of
what stands fast for it at the given time. It must then be left to closer testing
what theoretical and practical significance they may have; but practical
significance they can scarcely have if they have no theoretical significance
whatever. As regards the practical significance itself, finally, Lange gives no
definite criterion, no measure. Where he indicates his ethical conception, he
says that the natural sympathy which for Auguste Comte was the foundation of
ethics is not enough for an idealist ethics, and he regards the idea of a
totality to which we belong as the a priori foundation of ethics, by which alone
the moral law's inexorable strictness becomes intelligible. But at once he
becomes doubtful again: does the idea not often lead astray, and may it not be
asked, especially with respect to the religious systems, whether it is not better
simply to give oneself over to the ennobling effect of natural sympathy than to
listen to prophetic voices which only too often have led to the most terrible
fanaticism? --- A closer investigation of the relation between formalism and
realism in ethics Lange has not given. In his ethical ideas he was a disciple of
Kant, or rather of Schiller, though he saw the limitation of their conception.

His doctrine of the symbolic character and practical value of the religious ideas
Lange applies not only to philosophical speculation but also to the positive
religions. When materialism attacks Christianity as charcoal-burner's faith, it
overlooks --- Lange maintains --- the possibility of a freer, ideal conception of
the content of the ecclesiastical traditions. One should look to the ideas' value
and not to their literal formulation. In a certain sense religion's ideas are
imperishable. Who will refute a mass of Palestrina's, or establish that Raphael's
Madonna is an error? In the inwardly moved spiritual life of the genuinely devout
something stirs which the ideally minded freethinker must understand and must
sympathize with. Lange desires a reinterpretation of the ecclesiastical
doctrines, in the spirit of Kant, Fichte and Schleiermacher, so that
Christianity's ethical and ideal content may be drawn out of its literal husk;
Hegel's philosophy of religion was for him too conservative. Direct polemic
against positive religion he does not want, lest the people's unity and the
strength of its spiritual life should suffer. Of the two, he would rather
sacrifice enlightenment for a time than strength, provided only that freedom of
teaching is secured and priestly rule kept away. He has an open eye for the great
oppositions of our time, and expects from philosophy, as a contribution to
softening the struggle, above all a more penetrating understanding of the nature
and the laws of human development. He closes his brilliant work by expressing the
hope that serious work in this direction will bear its fruit. ---

The great interest in Kant's philosophy, to which Lange's work partly testified
and which it partly awakened, led to an eager study of the great philosopher whom
Romanticism's philosophy believed it had overcome. On several points disagreement
showed itself as to how Kant was to be understood. A rich Kant literature
blossomed, which in several respects led to a more exact conception of Kant, his
task, his method and his position in the history of philosophy. Even if one often
lost oneself in investigations of more philological than philosophical
significance, it was nevertheless of great importance that Kant should be
understood more correctly than before. This is due especially to works by
\textsc{Cohen}, \textsc{Fr.~Paulsen}, \textsc{Laas} and \textsc{Vaihinger}, to
which were later added the contributions to the illumination of Kant's
development and thought which \textsc{B.~Erdmann}, \textsc{R.~Reicke} and others
brought forward from Kant's unpublished notes and letters. In the great edition
of Kant's works begun by the Berlin Academy, the interest in critical philosophy
will set itself a fine memorial.

In a spirited paper \textit{Was uns Kant sein kann} (Vierteljahrsschr.\ für
wissensch.\ Philos.\ 1881) \textsc{Friedrich Paulsen} has sought to determine the
significance Kant can have for our time. That significance rests not upon a
direct resumption of his doctrine, but upon a treatment of the problems in his
spirit. So-called \emph{neo-Kantianism} therefore designates no closed school,
but an endeavour towards epistemological testing of the concepts we operate with.
\textsc{Wilhelm Wundt}, in his paper \textit{Was uns Kant nicht sein?}
(Philosophische Studien VII), has spoken against some neo-Kantians' unduly
dogmatic adherence to the old master, and has shown why Kant's doctrine cannot be
set up as an epistemology valid for the science of our own time.

\subsection*{b.\quad Eugen Dühring}

The year before Lange's ``Geschichte des Materialismus'' there appeared one of
the most excellent epistemological works of the century's last half, namely Eugen
Dühring's \textit{Natürliche Dialektik}. It takes its point of departure in
critical philosophy in so far as it means, in a more radical way, to take up
again Kant's problem, which Romanticism's ``deluge'' had washed away. Among later
German works it mentions only Trendelenburg's ``Logische Untersuchungen'' with
great recognition as a predecessor; among foreign thinkers Comte and Stuart Mill
are named, who have not, however, specially treated the book's subject. The book
is already now a literary rarity, since the author has not republished it ---
perhaps because his standpoint has since altered in a more positivist direction.
And yet it is surely the best of all Dühring's books, in form and in content
alike. It is an interesting contribution to the illumination of the relation
between critical philosophy and positivism, two tendencies which contend not only
in Dühring but in the whole of the most recent philosophy. By two fundamental
thoughts Dühring here already approaches positivism. Still more sharply than Kant
and Schopenhauer he maintains that the law of grounding --- the principle of
sufficient reason --- is only a law for our thinking, not for reality, which is
more comprehensive than our thinking. And he sets up, on the other hand, the
principle of the lacking ground, whereby the burden of proof is laid upon whoever
puts forward a new claim, what has hitherto been given being retained so long as
ground is lacking for abandoning it. By the first fundamental thought he rejects
rationalist dogmatism, by the second idealism and dualism; through both he works
his way forward to building upon the ground of the given reality, without the
admixture of representations not drawn from it. He continues Ludwig Feuerbach's
work and is positivism's most important representative in Germany. He himself
calls his philosophy the philosophy of reality. The great fundamental features of
reality, which experience teaches us to know, are to form the basis both of a
view of life and of the conduct of life. Philosophy is for him not merely theory,
but the expression of a personal disposition.

In Dühring's view the nineteenth century has a decidedly reactionary character
--- the technical field alone forms an exception. The seventeenth century stood
for him as the classical age in scientific respects. Serious science begins only
with it, and our century has brought forth nothing that can eclipse it. Dühring
is thinking here above all of Galilei, Huyghens and Newton in natural science,
and of Bruno, Hobbes and Spinoza in philosophy. It is the fundamental thoughts of
these great investigators that we still live upon. It is characteristic of
Dühring that the great, vigorous and fruitful initiating thoughts have for him
far greater significance than the special working-out and the confirmation by
experience. This hangs together with his historical eye and with his sense for
what is inspired; but it leads him to misjudge the importance of the great
thoughts finding their confirmation in particular experiences --- an importance
which a positivist of all people ought least to misjudge. This tendency in
Dühring, which was surely nourished by his blindness preventing him from seeing
and orienting himself for himself, increased in the course of his development and
became of fateful significance not only for his philosophy but also for his
outward circumstances.

While he praises the seventeenth century as the age of great thoughts, he praises
the eighteenth for the reforming application of those thoughts. When his own time
seems to him poor, he turns with all the greater hope towards the future. He
feels himself the herald of a tendency in thought and in life which will be able
to win currency only in a distant future. His great energy and fullness of
thought are united with the reformer's self-feeling --- unfortunately also with a
bitterness and a suspiciousness which his fate has nourished, but which make his
later writings in many parts unpleasant through coarse attacks upon those he
calls his enemies. His writings have a strongly personal stamp, in the word's
good and in its less good sense alike. He has himself depicted his personality
and his life in the remarkable book \textit{Sache, Leben und Feinde. Als
Hauptwerk und Schlüssel zu seinen sämmtlichen Schriften} (Karlsruhe und Leipzig
1882). He is right that his cause and his life hang closely together, and to his
honour it must be said that with him the cause stands before the life; less right
is he perhaps in supposing that his enemies hang together with his cause and his
life as closely as he suggests in the title.

\textsc{Eugen Dühring} was born in Berlin in 1833. His family came --- a thing he
set great store by --- from Sweden. By his father he was brought up in Rousseau's
spirit, under the influence of a liberal religious way of thinking. The
upbringing developed independence and firmness in his thinking. Mathematics and
astronomy early became principal interests. After his father's death he entered a
Gymnasium which was also a boarding house, but where he was as little satisfied
with the spiritual as with the bodily fare. He then studied law and worked for a
time as a practising lawyer. An eye disease which brought blindness with it made
this activity impossible, and gave occasion for his resolving to use his great
store of knowledge and thought to enter upon a scientific writer's career. His
wife, and later his son, was his secretary. A series of philosophical, economic
and scientific-historical works testify that the energetic thinker did not put
his trust in vain in his inner resources when the outer ones were closed to him.
Of his writings there are to be named, besides the one already mentioned,
\textit{Der Werth des Lebens} (1865), \textit{Cursus der Philosophie} (1875) ---
reworked as \textit{Wirklichkeitsphilosophie} (1895) --- \textit{Logik und
Wissenschaftstheorie} (1878) and \textit{Ersatz der Religion durch
Vollkommeneres} (1883). He worked as a Privatdocent at the University of Berlin,
where his lectures were among the best attended. His indignation at having been
passed over in the filling of professorships, and his criticism of conditions at
the university, brought him into sharp opposition to the leading circles there;
and when he came to believe that Helmholtz had usurped the credit for the
discovery of the law of the conservation of energy, which belonged to Robert
Mayer --- in whom Dühring saw the Galilei of the nineteenth century --- his
indignation found particular vent in repeated violent attacks upon the famous
physicist. The ministry, on the recommendation of the philosophical faculty,
punished him by depriving him of the right to lecture. That this was an unfair
harshness towards the blind philosopher seems, from everything available, to be
clear, even if his polemic was improper in content and in tone alike. The charge
Dühring brought against Helmholtz becomes intelligible only through his not
having followed the literature carefully, on account of his blindness; for
Helmholtz had in fact, as soon as he discovered that he had a predecessor in
Mayer, neglected no opportunity of giving him the honour that was his due. But
there operated here in Dühring as well that overvaluation of first ideas as
against their exact carrying through which we have mentioned above; what let him
place Robert Mayer so high above all his competitors was the same fundamental
consideration that let him set the seventeenth century so far above the
nineteenth. And finally his suspiciousness contributed, probably nourished by the
isolation blindness brought with it, and which let him find ``enemies'' on every
side. Next after the professors, the Social Democrats and the Jews were his
``enemies''. ``Opponents'' concerned with truth and not with their own advantage
he does not think he has met upon his way. (Leben, Sache und Feinde, Kap.~14.) He
is so certain of having all good powers on his side that he says at the close of
his autobiography: ``\textit{Ueberhaupt sind meine Feinde wesentlich keine andern
Elemente gewesen, als diejenigen, welche auch Feinde der besseren Menscheit
sind!}''

After working for some years by giving free lectures, Dühring has now withdrawn
to a small town near Berlin, occupied with natural-scientific and
literary-historical labours. A great and noble force has here been pushed aside,
though not broken, by misfortune and by his own fault and others'.

\subsubsection*{α.\quad Epistemology}

In his investigation of knowledge Dühring works (especially in ``Natürliche
Dialektik'') entirely in the spirit of critical philosophy. He seeks to exhibit
our knowledge's peculiarity in form and in procedure, in order to be able to
decide how far we may regard its forms and its results as expressions of
existence. He thus attempts a critical testing of the relation between thinking
and reality.

A principal point for him here is that thinking always seeks \emph{continuous
connection} and always seeks to \emph{continue} its course. In mathematics the
concept of infinity arises thereby, since no limits can be set to a magnitude's
decrease or increase. In logic we enter, by virtue of the principle of sufficient
reason, into an infinite series, one Why constantly succeeding another. Against
this tendency to continual connection and continuation Dühring sets what he calls
\emph{the law of the determinate number}. Every \emph{actual} given, every
\emph{factual} result is limited, and exhibits a determinate number of parts in
time and in space. Infinity and unlimitedness mean only the possibility of
continuation --- a possibility not always really warranted. And where the
continuation actually comes, it happens through successive additions of single,
determinate elements. --- There appears here, then, according to Dühring, a
characteristic difference between thought and reality.

Thought must be supple and untiring in order to be able to follow reality in all
its oscillations, in order to be able to rise and sink with the phenomena and to
discover all relations of dependence between them. But the law it must follow in
order to fulfil its task must not --- Dühring insists --- be transferred to
nature, to actual existence itself. This is seen clearly by comparing the real
science with the formal. In astronomy the determinate given magnitudes stand in
contrast to pure mathematics' unlimited possibility of production. In chemistry
the atomic numbers are examples of the law of the determinate number. Indeed, all
really given numbers and magnitudes are finite: there is a determinate number of
globes in the heavens, a determinate number of atoms in matter; at every moment a
wholly determinate number of parts of time has elapsed; the earth has moved round
the sun a certain determinate number of times; the causal series consists of a
finite number of members.

With great resoluteness Dühring draws from the law of the determinate number the
conclusion that the process of nature --- the sum of all changes in nature ---
must have had a beginning. An infinite regress cannot be thought. Prior to the
time filled with changes there lies, then, an eternal being, in which the
differences, the inner distinctions that bring succession and change with them,
had not yet appeared \emph{as successive}. Simultaneous multiplicity may very
well have been present, and must have been present, for it contains the germs
--- or, as Dühring says, the roots --- of all that we call process and apprehend
as determinate existence (Wirklichkeitsphilosophie, p.~15). Time and the causal
series first came to be with the transition from that homogeneous and
unchangeable condition to the unfolding of differences and changes. --- This
regressive conclusion of the causal series in an unchangeable condition is for
Dühring satisfying not only from the standpoint of thought but also from that of
feeling.

But Dühring has emphasized the opposition between thought's ideal continuity and
the reality of experience --- always given in separation and parcelling --- so
strongly that he ends in an insoluble riddle. For how is a transition from
absolute equilibrium and homogeneity to succession and change possible? He has
here, in his positivist thinking, come to face the same problem that Böhme and
Schelling had put forward in mythologizing and speculative form. He would soften
the paradox by pointing out that every time a new phenomenon [a new quality]
arises within our experience, we really have an absolute beginning. (Cursus p.~79.
Logik p.~191.) But the difference is that what is new and beginning in our
experience awakens wonder and gives thought a problem that can occupy it, whereas
that sharp opposition between absolute identity and change makes thought's work
hopeless once for all. Not merely a difference but a conflict between thinking
and reality is what Dühring here comes to lay down. --- When one has the choice
between assuming such a conflict and assuming that existence is as incapable of
being concluded as thinking is, this latter possibility seems preferable. There
is a constant interplay between thought and reality: as inquiry advances there
shows itself new continuity behind the apparent discontinuity, and so far reality
confirms thought's ideal; but new discontinuity is also discovered where there
seemed to be continuity in what is given, and the problem of how far reality can
be understood is then posed anew. In Dühring the dogmatist has got the upper hand
of the critical philosopher, on account of his great zeal to let reality come
into its own.

On the other hand Dühring lays very great weight upon the inner connection
between thinking and reality. Existence itself continues on into our thinking,
where this unfolds according to its true essence. To the relation between
premisses and conclusion (ground and consequent) there corresponds, in the real
processes, the relation between cause and effect. To logic's principle of
identity there corresponds the single thing's identity with itself,
notwithstanding that it occurs in different connections. To our thinking's
combinations and deductions there corresponds nature's unity and the interplay of
forces in the arising of phenomena. We do not thereby conceive nature on the
analogy of human consciousness, but as operating in such a way that human
knowledge gets a material to investigate that corresponds to its own
constitution. And even if our analysis will never be able to exhaust the great
connection in nature, the necessity does not cease of maintaining the kinship
between what operates in things and what operates in the understanding.

What appears in our scientific principles is according to Dühring the fruit of
forces that have long worked in a determinate direction before coming to clear
consciousness. Consciousness and thought rest upon something that is as little
consciousness and thought as the moving mechanical force is itself a phenomenon
of movement. But it is a warranted hypothesis to represent to ourselves the
ground of our thinking activity --- a ground not appearing in consciousness ---
as agreeing with the ground of nature's connection showing itself to us to be
ordered in accordance with the understanding. (Nat.\ Dial.\ p.~155, 225. Logik
p.~173. Leben, Sache etc.\ p.~303.) What makes us think is thus the same as what
makes nature operate.

This last train of thought now gradually gains such preponderance over Dühring
that in his later writings he conducts a violent polemic against the critical
philosophy, which maintains that we do not know things in themselves, because our
knowledge is determined not merely by the nature of things but also by our own
nature. He sees in that conception a bogging-down of science, a hypocritical
accommodation of thought's results to what the need to assume a world beyond
demands. Here again it is the law of the determinate number that underlies the
matter: number is an example of our commanding concepts which express what is, in
such a way that between reality and the concept no insertion conditioned by the
constitution of the I takes place at all. (Logik p.~165.)

As sharply as Dühring elsewhere emphasizes the opposition between thinking and
being, so definitely does he here emphasize the possibility of their completely
covering one another. There he has formal knowledge chiefly in view, here real
knowledge. But they cannot be kept wholly apart, as is shown by the fact that the
law of the determinate numbers precisely sets inquiry new tasks --- namely, to
find out why changes occur just when those determinate numbers are reached.

It is Dühring's merit with respect to recent epistemology that he sets the
problem of the relation between quality and continuity, between the determinate
turning points and the connected processes, in so clear and distinct a light. His
whole standpoint is determined by a contention between criticism and positivism,
and contains an invitation to a closer investigation of the relation between
those two tendencies.

\subsubsection*{β.\quad Conception of the World}

Dühring lays great weight upon the development of a \emph{world-schematism},
which is to give a fully formed picture of actual existence's general features.
In this the task of the philosophy of reality consists. It is to lead to a system
of experience, whereas the earlier metaphysics was a kind of alchemy which sought
by mystical ways to fathom existence's essence. --- At this point he recalls
Comte and Spencer. He is in his ``schematism'' more abstract than Spencer and
less original than Comte.

My system, he says, knows not two realities but only a single one --- nature and
its parts. He approaches materialism in emphasizing that everything in the world
has a material side, and that by following the material phenomena in their
intuitable connection one must be able to find everything that has reality. But
he will not stop at the mere phenomena. Thereby he distances himself from ``the
usual, rather coarse materialism''. He often uses the word \emph{matter} in the
sense of \emph{being, world-medium, the bearer of all that is actual}. But he
expressly emphasizes that the concept of matter must be widened and deepened so
that it becomes clear that the dispositions to the phenomena of consciousness are
found in nature just as much as the possibilities of all other phenomena.
\emph{The philosophical concept of matter} contains more than the physicists' and
the chemists', since it is to be the concept of what underlies all bodily and
spiritual states and comprehends the whole and full reality. (Cursus p.~62, 75.
Leben, Sache etc.\ p.~302\,f.) Materialism is for Dühring only a ``pedestal'', a
real foundation, an outer, mechanical frame for the conception of the world and
the view of life --- not the exhaustive, ample explanation of existence itself.

A doctrine one does not expect to find in Dühring is the assumption of
\emph{purposes} in nature. To understand him rightly on this point one must
notice that he distinguishes between purpose (\textit{Zweck}) and intention
(\textit{Absicht}). (Logik p.~203.) When he ascribes validity to the concept of
purpose for our conception of nature, it is because the forces and dispositions
operating in nature always exhibit certain tendencies, work in certain
directions, and above all because they \emph{operate together} and thereby make
nature appear as a systematic unity. Those tendencies do not always lead --- as
little in outer nature as in conscious human life --- to completed results. But
whether the goal is reached or not, \emph{the combined operation of manifold
elements} appears as a characteristic principal feature of all processes. In
apprehending nature one must attend not merely to the general laws but also to
the directions and tendencies of the processes that go on, and to the types and
forms in which nature's activity appears. Nature forms a scale of stages, the
lower forms of existence being the foundation of the higher. The measure by which
we distinguish between lower and higher forms can be sought only in the results
of the different tendencies, in the degree to which they lead to their natural
goal. And the last goal must be the arising of beings which not only exist and
operate but also become conscious of their existence and their operating: only as
the foundation of the history of conscious beings has the history of the
mechanical whole of nature a meaning; a thoroughly unconscious world would ---
says Dühring (Cursus p.~104) --- be a foolish half-measure, a theatre without
actors and without spectators. And that the phenomena of consciousness are, on
the other hand, essentially to be regarded as purposes is seen from the fact that
unconscious and conscious operation can accomplish one and the same thing; it
would be mere luxury for men and animals to become conscious of their impulses
and motives, if it were only a question of attaining outer results, and if we
were not here facing something that was in itself purpose. Our feelings in
particular --- the feeling of love, the social feeling --- do not have their
value merely by being means or motives leading to actions necessary for life's
subsistence; they have immediate value as expressions of a higher stage of life
than the purely elementary one. (Cursus p.~158\,f. Der Werth des Lebens,
2.~Aufl., p.~241.)

If one were to object to this last thought that conscious life is nevertheless
not always bound up with pleasure, but that pain occupies so great a place that
conscious life does not seem to be a real ``purpose'', Dühring would point to
what he calls \emph{the law of difference} --- a law which under different forms
and names plays a great part in recent psychology. (The most suitable name would
be: the law of ratio.) For Dühring it is not merely a psychological law but a
world-law. The antagonism of forces is a principal feature in the
world-schematism. All expression of force, all movement and all development is
conditioned by distance, difference, opposition. All consciousness too
presupposes difference. Without difference, no sensation. Within consciousness
the same holds. Life's play would lose its attraction if there were no resistance
and hindrance to overcome, and if need and satisfaction did not succeed one
another. It is existence's rhythms and differences that make it valuable for us:
our feeling of life is set in motion by the transitions between mutually
different states. Without what is harsh, bitter and painful, the deep
satisfaction in life would not be felt. --- The same emphasis upon feeling's
being determined by oppositions which made Schopenhauer a pessimist makes Dühring
an optimist. Schopenhauer is, as regards the ``law of difference'', more
consistent than Dühring, in so far as he assumes that unrest and movement lie as
an original tendency in the world-ground's nature; whereas for Dühring's
philosophy it becomes a great riddle how successive differences and rhythmically
changing states could have succeeded the resting being which that philosophy
assumes in virtue of the law of the determinate number.

The \emph{combination} of the different forces is what makes development
possible. Whether this combination leads to a state of equilibrium or to a new
movement, the tendencies of the several forces are preserved in the development's
result. A constant strife between the forces would be an absurdity --- but it is
not found in nature. The avoidance of the absurd is a natural-logical principle.
Only in exceptional cases does this principle lead to untenable forms being
annihilated and dissolved. As a rule nature takes the positive way. Therefore, on
Dühring's view, the idea of the struggle for existence is false: it one-sidedly
emphasizes the negative point of view and lays the weight upon antagonism instead
of upon combination. Dühring gives Lamarck's theory the preference over Darwin's.

\subsubsection*{γ.\quad Ethics}

The ethical has according to Dühring its foundation in the human impulses. Only
its foundation --- for nature can err, and therefore a closer development and
correction is necessary. The higher a being stands, the greater the possibility
that it goes wrong. Besides, what matters ethically is to combine elements which
from nature's hand are not yet connected, but are perhaps even in mutual strife.

% [p. 565]
The germs of the good Dühring finds, like Comte, in the sympathetic instincts.
Nature has itself seen to it that another's suffering determines our own feeling
in a painful way, and there is perhaps no feeling that has undergone so
conspicuous a development in the course of culture as this one. It does not lead
to subduing and abolishing our own individuality. On the contrary, ethical
development consists both in an individualization and in a socialization; the two
necessarily belong together, since the full development of the single
individuality first becomes possible in a more highly developed society. Under
more imperfect social conditions the free, distinctive development of individuals
is hindered. But since our present states are essentially built upon compulsion,
they constantly hamper that development. A free society --- the germs of which
are given in the organizations formed by free association --- will first let both
individuality and community come into their own. In such a society all relations
of production and of consumption would be brought under social direction, so that
the individual's whole interest can be attached to the work and not to the work's
yield. Only thereby does a thoroughgoing ennoblement of life become possible.
From socialism Dühring's conception differs partly in that he thinks of free
labour communes and not of an absolute centralization of
production,\footnote{On Dühring's position on the social question, and on the
different standpoints he has taken here, I refer to B. Friedländer: \textit{Die
vier Hauptrichtungen der modernen socialen Bewegung.} Berlin 1901.} partly in
that he does not regard it as necessary that development should be interrupted by
a historical miracle: the building of the future will on his view arise, not ---
as Karl Marx holds --- through evil increasing, but through the good growing
quietly. On this point too Dühring is decidedly an optimist, however much
baseness he may think he can establish in the world. He is a definite opponent of
pessimism. Only the pessimism of indignation (\textit{Entrüstungspessimismus}),
as represented by Byron and called forth by determinate social conditions, has
his approval; whereas he despises the Romantic ``pessimism of the beyond'' and
the ``pessimism of rottenness'' frequent in the higher classes.

Beyond the feeling for the good forces that work their way forward within the
human world, there arises a \emph{universal affect} when the gaze is raised to
the great unity of which the human world is only a single ramification --- when
the thinker feels himself as belonging to that unity and as moved by its forces
--- and when the thought of his own fate is taken up into the thought of the
great order of things, within which so many good germs have been able to unfold.
---

Not only by the problems it poses, but also by the noble, antique style in which
it is carried through, and by the connection it establishes between thinking and
personality, Dühring's philosophy stands as one of the most characteristic
attempts of thought in the generation that followed Romanticism's fall and
natural science's new rise at the middle of the nineteenth century. The year 1880
forms here a characteristic turning point, in that a circle of the more
considerable investigators conclude their activity about that time --- namely,
besides Dühring, Lotze, Fechner, Darwin and Spencer; shortly before, Lange had
succumbed. In a separate work some principal forms of the history of philosophy
in the following quarter-century are to be presented.

\end{document}
